\documentclass[11pt]{article}
% Shared preamble for manuscripts-v2 papers. Load with % Shared preamble for manuscripts-v2 papers. Load with % Shared preamble for manuscripts-v2 papers. Load with \input{../common/preamble.tex}
% immediately after \documentclass[11pt]{article}.
\ifdefined\pdftrailerid\pdftrailerid{}\fi
\usepackage[T1]{fontenc}
\usepackage{lmodern}
\usepackage[letterpaper,margin=1.02in]{geometry}
\usepackage{amsmath,amssymb,amsthm,mathtools,booktabs,array,microtype,enumitem}
\usepackage{tikz}
\usepackage{xurl}
\usepackage[colorlinks=true,allcolors=bluegreen]{hyperref}
\definecolor{bluegreen}{RGB}{22,83,91}
\usepackage[capitalize,nameinlink]{cleveref}

\theoremstyle{plain}
\newtheorem{theorem}{Theorem}[section]
\newtheorem{lemma}[theorem]{Lemma}
\newtheorem{proposition}[theorem]{Proposition}
\newtheorem{corollary}[theorem]{Corollary}
\newtheorem{conjecture}[theorem]{Conjecture}
\newtheorem*{theorem*}{Theorem}
\newtheorem{maintheorem}{Theorem}
\renewcommand{\themaintheorem}{\Alph{maintheorem}}
\theoremstyle{definition}
\newtheorem{definition}[theorem]{Definition}
\newtheorem{example}[theorem]{Example}
\newtheorem{question}[theorem]{Question}
\theoremstyle{remark}
\newtheorem{remark}[theorem]{Remark}
\crefname{maintheorem}{Theorem}{Theorems}
\crefname{conjecture}{Conjecture}{Conjectures}
\crefname{question}{Question}{Questions}
\crefname{example}{Example}{Examples}

\newcommand{\R}{\mathbb R}
\newcommand{\Q}{\mathbb Q}
\newcommand{\Z}{\mathbb Z}
\newcommand{\N}{\mathbb N}
\newcommand{\T}{\mathbb T}
\newcommand{\C}{\mathbb C}
\DeclareMathOperator{\den}{den}
\newcommand{\e}{\mathrm e}

% Immutable Lean source links at research-checkpoint-18.
\newcommand{\leancommit}{e4809cf561573aef588e6ffc491a35cdb79fd9ec}
\newcommand{\leanbase}{https://github.com/eldernos/MathematicsResearch/blob/\leancommit}
% \lean{File}{line}{label}: link to NormalityChecker/File.lean#Lline at the pinned commit.
\newcommand{\lean}[3]{\href{\leanbase/NormalityChecker/#1.lean\#L#2}{#3}}
% \repo{path}{label}: link to any repository path at the pinned commit.
\newcommand{\repo}[2]{\href{https://github.com/eldernos/MathematicsResearch/tree/\leancommit/#1}{#2}}

% Evidence labels (see papers/WAVE2-CONTRACT.md).
\newcommand{\evidence}[1]{\textup{[\emph{#1}]}}

\setlength{\emergencystretch}{1.5em}
\raggedbottom

% immediately after \documentclass[11pt]{article}.
\ifdefined\pdftrailerid\pdftrailerid{}\fi
\usepackage[T1]{fontenc}
\usepackage{lmodern}
\usepackage[letterpaper,margin=1.02in]{geometry}
\usepackage{amsmath,amssymb,amsthm,mathtools,booktabs,array,microtype,enumitem}
\usepackage{tikz}
\usepackage{xurl}
\usepackage[colorlinks=true,allcolors=bluegreen]{hyperref}
\definecolor{bluegreen}{RGB}{22,83,91}
\usepackage[capitalize,nameinlink]{cleveref}

\theoremstyle{plain}
\newtheorem{theorem}{Theorem}[section]
\newtheorem{lemma}[theorem]{Lemma}
\newtheorem{proposition}[theorem]{Proposition}
\newtheorem{corollary}[theorem]{Corollary}
\newtheorem{conjecture}[theorem]{Conjecture}
\newtheorem*{theorem*}{Theorem}
\newtheorem{maintheorem}{Theorem}
\renewcommand{\themaintheorem}{\Alph{maintheorem}}
\theoremstyle{definition}
\newtheorem{definition}[theorem]{Definition}
\newtheorem{example}[theorem]{Example}
\newtheorem{question}[theorem]{Question}
\theoremstyle{remark}
\newtheorem{remark}[theorem]{Remark}
\crefname{maintheorem}{Theorem}{Theorems}
\crefname{conjecture}{Conjecture}{Conjectures}
\crefname{question}{Question}{Questions}
\crefname{example}{Example}{Examples}

\newcommand{\R}{\mathbb R}
\newcommand{\Q}{\mathbb Q}
\newcommand{\Z}{\mathbb Z}
\newcommand{\N}{\mathbb N}
\newcommand{\T}{\mathbb T}
\newcommand{\C}{\mathbb C}
\DeclareMathOperator{\den}{den}
\newcommand{\e}{\mathrm e}

% Immutable Lean source links at research-checkpoint-18.
\newcommand{\leancommit}{e4809cf561573aef588e6ffc491a35cdb79fd9ec}
\newcommand{\leanbase}{https://github.com/eldernos/MathematicsResearch/blob/\leancommit}
% \lean{File}{line}{label}: link to NormalityChecker/File.lean#Lline at the pinned commit.
\newcommand{\lean}[3]{\href{\leanbase/NormalityChecker/#1.lean\#L#2}{#3}}
% \repo{path}{label}: link to any repository path at the pinned commit.
\newcommand{\repo}[2]{\href{https://github.com/eldernos/MathematicsResearch/tree/\leancommit/#1}{#2}}

% Evidence labels (see papers/WAVE2-CONTRACT.md).
\newcommand{\evidence}[1]{\textup{[\emph{#1}]}}

\setlength{\emergencystretch}{1.5em}
\raggedbottom

% immediately after \documentclass[11pt]{article}.
\ifdefined\pdftrailerid\pdftrailerid{}\fi
\usepackage[T1]{fontenc}
\usepackage{lmodern}
\usepackage[letterpaper,margin=1.02in]{geometry}
\usepackage{amsmath,amssymb,amsthm,mathtools,booktabs,array,microtype,enumitem}
\usepackage{tikz}
\usepackage{xurl}
\usepackage[colorlinks=true,allcolors=bluegreen]{hyperref}
\definecolor{bluegreen}{RGB}{22,83,91}
\usepackage[capitalize,nameinlink]{cleveref}

\theoremstyle{plain}
\newtheorem{theorem}{Theorem}[section]
\newtheorem{lemma}[theorem]{Lemma}
\newtheorem{proposition}[theorem]{Proposition}
\newtheorem{corollary}[theorem]{Corollary}
\newtheorem{conjecture}[theorem]{Conjecture}
\newtheorem*{theorem*}{Theorem}
\newtheorem{maintheorem}{Theorem}
\renewcommand{\themaintheorem}{\Alph{maintheorem}}
\theoremstyle{definition}
\newtheorem{definition}[theorem]{Definition}
\newtheorem{example}[theorem]{Example}
\newtheorem{question}[theorem]{Question}
\theoremstyle{remark}
\newtheorem{remark}[theorem]{Remark}
\crefname{maintheorem}{Theorem}{Theorems}
\crefname{conjecture}{Conjecture}{Conjectures}
\crefname{question}{Question}{Questions}
\crefname{example}{Example}{Examples}

\newcommand{\R}{\mathbb R}
\newcommand{\Q}{\mathbb Q}
\newcommand{\Z}{\mathbb Z}
\newcommand{\N}{\mathbb N}
\newcommand{\T}{\mathbb T}
\newcommand{\C}{\mathbb C}
\DeclareMathOperator{\den}{den}
\newcommand{\e}{\mathrm e}

% Immutable Lean source links at research-checkpoint-18.
\newcommand{\leancommit}{e4809cf561573aef588e6ffc491a35cdb79fd9ec}
\newcommand{\leanbase}{https://github.com/eldernos/MathematicsResearch/blob/\leancommit}
% \lean{File}{line}{label}: link to NormalityChecker/File.lean#Lline at the pinned commit.
\newcommand{\lean}[3]{\href{\leanbase/NormalityChecker/#1.lean\#L#2}{#3}}
% \repo{path}{label}: link to any repository path at the pinned commit.
\newcommand{\repo}[2]{\href{https://github.com/eldernos/MathematicsResearch/tree/\leancommit/#1}{#2}}

% Evidence labels (see papers/WAVE2-CONTRACT.md).
\newcommand{\evidence}[1]{\textup{[\emph{#1}]}}

\setlength{\emergencystretch}{1.5em}
\raggedbottom

% Author information approved by the repository owner on 2026-11-25:
%   name: Christopher Anderson; designation: Independent Researcher.
% No academic affiliation, ORCID or personal contact details have been supplied.
% Do not add any that the owner has not approved.
\newcommand{\OwnerAuthor}{Christopher Anderson}
\newcommand{\OwnerAffiliation}{Independent Researcher}
\newcommand{\OwnerContact}{Corrections and comments: \url{https://github.com/eldernos/MathematicsResearch/issues}}
\newcommand{\ReleaseLabel}{manuscripts-v2.0.0}
\newcommand{\LicenseLine}{\copyright\ 2026 Christopher Anderson. Licensed under \href{https://creativecommons.org/licenses/by/4.0/}{CC BY 4.0}.}
% \makebyline{version}{date}: centred author block used under the title.
\newcommand{\makebyline}[2]{%
  {\large \OwnerAuthor\par}%
  \ifx\OwnerAffiliation\empty\else{\small \OwnerAffiliation\par}\fi
  \vspace{0.3em}
  {\small Version #1 \enspace$\cdot$\enspace #2 \enspace$\cdot$\enspace \ReleaseLabel\par}%
  {\footnotesize \OwnerContact\par}%
  {\footnotesize \LicenseLine\par}%
  \hypersetup{pdfauthor={\OwnerAuthor}}}

\hypersetup{pdftitle={Synthetic Arithmetic Expansions and Prescribed Discrepancy}}
\newcommand{\dist}[1]{\lVert #1\rVert}
\newcommand{\fr}[1]{\{#1\}}
\newcommand{\Prob}{\mathcal P}
\newcommand{\Minv}{\mathcal M_T}
\DeclareMathOperator{\ord}{ord}
\DeclareMathOperator{\lcm}{lcm}
\newcommand{\ttt}[1]{\texttt{#1}}
\newcommand{\leanP}[3]{\leanat{\phaseXIXcommit}{#1}{#2}{#3}}
\setlength{\emergencystretch}{3em}

\begin{document}
\begin{center}
{\LARGE\bfseries Synthetic Arithmetic Expansions\\[3pt] and Prescribed Discrepancy\par}
\vspace{10pt}
\makebyline{1.0}{25 November 2026}
\end{center}

\begin{abstract}
Fix an integer base $b\ge2$, a prime $p\nmid b$, and let $t$ be the order of $b$ modulo $p$.
For an integer $d\ge2$, insert the coefficient $p^{-s v_p(k)}$ at base-$b$ position $tk^d$ for every
$k\ge1$. For $s=0$ the extreme discrepancy tends to one, while for every integer $s\ge1$ the number is
normal to base $b$; this switch holds on every strictly increasing schedule of multiples of $t$
dominating $k^2$, and it has been checked by the Lean~4 kernel together with the exact reduced
denominator of every shifted prefix. For odd $p$ and $d\ge s+2$ complete periods give an explicit power
rate, and for $\xi=\sum_{k}3^{-v_3(k)}2^{-2k^3}$ a short self-contained proof gives $D_N\le200N^{-1/3}$.
Every radix orbit is, after a finite shift, the orbit of a positive expansion of this kind, so
qualitative realization of discrepancy limits and limit sets follows from classical results. Choosing
the rational state after each cubic insertion from two explicit targets, with coefficient denominators
at most $pk$, realizes every limiting discrepancy in $[0,1]$ with error $O(N^{-1/3})$ uniformly in the
limit, and gives explicit arithmetic realizations of every pair of lower and upper limits and of every
nonempty compact connected set of invariant measures. For odd $p$, $s\ge1$ and an integer $d\ge3$, every
nonconstant rational polynomial of total degree at most two in the numbers of this family with distinct
orders $s$ on a common schedule is normal; in particular $\xi^2$ is normal and $(\xi,\xi^2)$ is jointly
normal. The proof shows that the limiting $p$-adic mass of the shifted prefixes has Lebesgue-null level
sets, by counting isolated jumps with van der Corput's method and a second-moment argument. Finally a
positive binary series using every odd prime as a denominator prime is normal and Liouville with
$D_N=O(1/\log N)$, while positive rational multiples of its jointly normal components can sum to any
positive rational.
\end{abstract}

\section{Introduction}\label{sec:intro}

Throughout, $b\ge2$ is an integer, $\T=\R/\Z$, $T(y)=by\bmod 1$, $\fr{y}$ is the fractional part
and $\dist{y}$ the distance to the nearest integer. For real $x$ and $N\ge1$ the \emph{extreme
discrepancy} of the first $N$ points of the radix orbit is
\[
D_N(x)=D_N(x;b)=\sup_{0\le a<c\le 1}\left|\frac{\#\{0\le n<N:\fr{b^nx}\in[a,c)\}}{N}-(c-a)\right|.
\]
A real number is normal to base $b$ if and only if $D_N(x)\to0$ \cite[Chapter~1, Theorem~8.1,
p.~70 and Chapter~2, Theorem~1.1, p.~89]{KN}. For a prime $p$ we write $v_p$ for the $p$-adic
valuation and $\ord_m(b)$ for the multiplicative order of $b$ modulo $m$.

The numbers studied here are \emph{synthetic expansions}: sums $\sum_k c_kb^{-n_k}$ with
positive rational coefficients placed at a sparse arithmetic schedule $n_k$, where the coefficients
are chosen so that the rational part $b^n\sum_{n_k\le n}c_kb^{-n_k}$ seen after each shift has a
prescribed reduced denominator or a prescribed residue. The first family is defined by a single
valuation weight.

\begin{maintheorem}[Normality switch]\label{thm:A}
Let $p\nmid b$ be prime, $t=\ord_p(b)$, $s\ge0$ an integer, and let $(n_k)_{k\ge1}$ be strictly
increasing positive multiples of $t$ with $n_k\ge k^2$. Put
\[
X_s=\sum_{k\ge1}\frac{1}{p^{s v_p(k)}\,b^{n_k}}.
\]
Then $D_N(X_0)\to1$, and $D_N(X_s)\to0$ for every $s\ge1$, as $N\to\infty$ through all integers.
\end{maintheorem}

\Cref{thm:A} is kernel-verified (\Cref{thm:switch}); so is the exact reduced denominator of every
shifted prefix (\Cref{thm:den}), which is the arithmetic mechanism: a positive order $s$ keeps a
pole of exact order $s\lfloor\log_pK\rfloor$ after the $K$th insertion, while $s=0$ retires all
denominators. We give a written proof of the following rate.

\begin{maintheorem}[Complete-period rate]\label{thm:B}
Let $p\nmid b$ be an odd prime, $t=\ord_p(b)$, and $s\ge1$, $d\ge s+2$ integers. For
$X=\sum_{k\ge1}p^{-sv_p(k)}b^{-tk^d}$ and $N\ge2$,
\[
D_N(X)\ll_{b,p,s,d} N^{-s/d}+N^{-(d-s-1)/d}+N^{-(d-1)/d}\log N .
\]
For the binary number $\xi=\sum_{k\ge1}3^{-v_3(k)}2^{-2k^3}$ one has $D_N(\xi)\le 200N^{-1/3}$
for every $N\ge1$.
\end{maintheorem}

The second half is proved first, in \Cref{sec:example}, by an argument that uses only the fact
that $2$ is a primitive root modulo every power of $3$. The general estimate is \Cref{thm:rate}.

Changing the coefficient after each insertion lets one prescribe the rational state, and hence the
orbit, on each gap. Retiring a prescribed weighted proportion of gaps gives every discrepancy limit
with an explicit error.

\begin{maintheorem}[Prescribed discrepancy]\label{thm:C}
Let $p\nmid b$ be an odd prime and $t=\ord_p(b)$. For every $\delta\in[0,1]$ there are rationals
$c_k\in(0,1]$ whose reduced denominators divide $p^{1+\lfloor\log_pk\rfloor}\le pk$ such that
$x_\delta=\sum_{k\ge1}c_kb^{-tk^3}$ satisfies
\[
|D_N(x_\delta)-\delta|\le C_{b,p}N^{-1/3}\qquad(N\ge1),
\]
with $C_{b,p}$ independent of $\delta$. On the same schedule $tk^3$ and with the same denominator bound,
every pair $0\le\alpha\le\beta\le1$ occurs as $(\liminf D_N,\limsup D_N)$, with
$D_N=\Lambda_N/N+O(N^{-1/3})$ for an explicit weighted count $\Lambda_N$. Every nonempty compact
connected set of $T$-invariant Borel probability measures on $\T$ occurs, again on the schedule $tk^3$,
as the set of weak limit points of the orbit measures of an expansion whose state after the $k$th
insertion is a purely periodic rational of period at most $k$.
\end{maintheorem}

Qualitatively, \Cref{thm:C} adds nothing to classical realization. Every radix orbit is, after
finitely many shifts, the orbit of a sum $\sum_kc_kb^{-n_k}$ with rationals $c_k\in(0,1)$ whose
denominators are powers of $p$, on any prescribed increasing schedule (\Cref{rem:qualitative}). The
realization of invariant-measure limit sets by digit sequences \cite{Colebrook,Sigmund74,DGS}, combined
with weak continuity of $D$ (\Cref{lem:Dcont}), therefore already yields every limit $\delta$, every pair
$(\alpha,\beta)$ and every limit set by such expansions. What \Cref{thm:C} adds, in the sources we
examined, is quantitative and arithmetic: an error $O(N^{-1/3})$ at every prefix, uniform in $\delta$;
coefficient denominators bounded by $pk$, with the target state chosen from two explicit values; and an
explicit arithmetic schedule of targets, given by the Beatty-type mask
$\lfloor k\delta\rfloor-\lfloor(k-1)\delta\rfloor$. The proofs are in \Cref{sec:program,sec:limitsets}.

The linear switch of \Cref{thm:A} rests on a single surviving pole. For squares that pole can
vanish: in the shifted prefixes of $\xi^2$ the top power of $3$ cancels on a set of shifts of positive
lower density, so the method of \Cref{thm:A} alone does not decide normality. A finer analytic argument
does.

\begin{maintheorem}[Quadratic images]\label{thm:D}
Let $p\nmid b$ be an odd prime, $t=\ord_p(b)$, $d\ge3$ an integer, and for integers $a\ge1$ put
$X_a=\sum_{i\ge1}p^{-av_p(i)}b^{-ti^d}$. For every finite set $S$ of positive integers, every nonconstant
polynomial $P\in\Q[z_a:a\in S]$ of total degree at most two gives a number $P((X_a)_{a\in S})$ normal
to base $b$. In particular every rational quadratic polynomial in $X_s$, $s\ge1$, is normal; the
distinct nonconstant monomials of degree at most two in the $X_a$ are jointly normal; $\xi^2$ is normal
to base $2$; and the binary shifts of $(\xi,\xi^2)$ are uniformly distributed in $\T^2$.
\end{maintheorem}

This is \Cref{thm:quadratic}, a written proof without a rate (\Cref{sec:quadratic}). The arithmetic
input is the exact all-prefix criterion of \Cref{prop:criterion}, which reduces normality to showing
that small reduced denominators occur on a set of shifts of density zero. For quadratic prefixes the
relevant quantity, after normalization, converges $p$-adically to a function $M(u)$ of the real
position $u$ of the shift in its $p$-adic shell, with values in $\Z_p$, whose reduction modulo each
power of $p$ is a finite step function. We show that every level set of $M$ is Lebesgue-null:
primitive jumps of $M$ on the grids $p^{-H}\Z$ are isolated from all competing jumps with positive
proportion, by an equidistribution estimate using the third-derivative test of van der Corput, and a
second-moment argument places almost every point in infinitely many isolated neighbourhoods.

Finally we test synthesis with infinitely many denominator primes.

\begin{maintheorem}[Infinite sums]\label{thm:E}
For an odd prime $\ell$ let $s_3=2$, $s_\ell=1$ for $\ell\ge5$, and
$Y_\ell=\sum_{j\ge1}\ell^{-s_\ell j}2^{-16^{\ell^j}}$. Then $Y=\sum_\ell Y_\ell$, summed over all odd
primes, is normal to base $2$ and a Liouville number, and $D_N(Y)=O(1/\log N)$. Every finite
subfamily of the $Y_\ell$ is jointly normal, and there are positive rationals $c_\ell$ with
$\sum_\ell c_\ell Y_\ell=1$.
\end{maintheorem}

Thus finite joint normality of positive terms does not survive infinite summation, while a
uniform arithmetic grid (\Cref{lem:grid}) at the prime $3$ survives the introduction of every other
odd prime (\Cref{sec:sums}).

\paragraph{Companion manuscript.} The perturbation estimates used here are finite comparisons
(\Cref{lem:compare}). Their sharp form, the identity $\sup_x|D_N(x+z)-D_N(x)|=1-D_N(z)$, together with
the classification of translations preserving discrepancy, the exact addition spectrum of $D_N$ and
the behaviour of $D_N$ under rational scalars, is the subject of the companion manuscript
\cite{CompanionH} in the same release. We quote the identity once (\Cref{thm:sharpcited}) and use it
only in \Cref{prop:tails}.

\paragraph{Relation to earlier work.} Arithmetic constructions of normal numbers by sparse series
go back to Stoneham \cite{Stoneham73} and Korobov \cite{Korobov72}; Kano \cite[Theorem~1]{Kano93} gives
a general criterion for series $\sum A_j/(a^{\lambda_j}b^{\mu_j})$, with strictly increasing exponents
and size conditions on the numerators, and records in its introduction Stoneham's constructions with
prescribed numerators. Bailey and Crandall \cite[Theorem~4.8]{BC02} prove normality for such series
with numerator one under growth conditions on the two exponent sequences, and Vandehey
\cite[Corollary~1.2 and Theorem~7.4]{Vandehey19} gives power-saving estimates for incomplete Korobov
sums on moduli with fixed prime support, together with a normality criterion for series
$\sum 1/(c_kb^{m_k})$ with numerator one in which $c_k\mid c_{k+1}$ strictly. Two denominator sequences
should be distinguished in \Cref{thm:A}. The insertion denominators $p^{sv_p(k)}$ return to $1$ whenever
$p\nmid k$, so the series is not term by term an instance of these criteria; the reduced denominators
$p^{s\lfloor\log_pK\rfloor}$ of the shifted prefixes do form a divisibility chain. Grouping consecutive
insertions into single terms does not help either: on the cubic schedule a group of two or more terms,
written over its last radix position, has size times that radix power tending to infinity, which is
incompatible with the bounded normalized numerators of these criteria (the test is recorded in the
ledger, \Cref{sec:formal}). We therefore do not apply a published series criterion. The positive side
of \Cref{thm:A} combines the exact prefix denominators (\Cref{thm:den}) with a fixed-support short-orbit
estimate (\Cref{lem:short}), the qualitative consequence of Vandehey's Corollary~1.2; the side $s=0$
is elementary, and the switch holds on every schedule dominating $k^2$. Rings of numbers normal to one
base and not to another were constructed by Wagner \cite{Wagner95}, an uncountable ring all of whose
nonzero elements are normal to base $b$ and not normal to base $ab$ for a prime $a\nmid b$ (as described
in \cite[\S1]{Kano93} and in the introduction of \cite{KS93}), and by Kano and Shiokawa
\cite[Theorem~5]{KS93} for $a$ not necessarily prime. Since nonzero polynomial expressions in the
generators lie in the ring, these are direct precedents for the normality of polynomial images of
suitable generators. Polynomial images of the standard Stoneham numbers are treated in \cite{PaperI},
and Mira Cult's public note \cite[Theorem~1.1]{MiraPoly} proves the normality of every nonconstant
rational polynomial in $\sum_kq^{-k}b^{-c^k}$ when $q>2$ and some prime divides $c$ and $q$ but not $b$
(see the comparison after \Cref{cor:flagship}). The ring generators use rapidly separated exponents,
and S08 uses the exponential schedule $c^k$. \Cref{thm:D} concerns specific generators on a polynomial
schedule, whose consecutive positions have ratio tending to one and whose weight exponents $sv_p(k)$
return to zero infinitely often; we have not verified that the generator hypotheses of the ring
constructions hold for them. Its conclusion is limited to total degree two, and it does not produce a
ring.
Normal Liouville numbers are known explicitly \cite{KS93,BC02}, absolutely normal Liouville numbers
exist \cite{Bugeaud02}, and a computable one was constructed in \cite{BHS15}; \Cref{cor:liouville} is a
further explicit example inside the switch family. For a family of signed Stoneham expansions the
attainable upper discrepancies form an explicit proper subset of $[0,1]$ with a gap above zero
\cite{PaperIII}, whereas \Cref{thm:C} shows that positive coefficients on a cubic schedule reach every
value with a uniform rate. Recent public work proves the normality of sums of two Stoneham numbers \cite{SudeSum} and the
joint normality of Stoneham series on finitely many distinct primes and orders \cite{MiraJoint};
\Cref{thm:combos} is a result of the same kind for the different series of \Cref{thm:E}. Normal numbers
with very small discrepancy are constructed in \cite{Levin99,ABSS17}; the rates of \Cref{thm:B,thm:C}
are much weaker and are included because they are elementary and hold for the specific arithmetic
families. We have not located \Cref{thm:A,thm:D,thm:E} in the sources examined; a detailed
comparison, including recent public repositories, is kept in the ledger accompanying the source files
(\Cref{sec:formal}).

\paragraph{Evidence labels.} Each numbered statement carries one label:
\emph{Kernel-verified} (the complete statement is a Lean declaration checked by the kernel, linked
at a pinned commit), \emph{Lean-supported paper proof} (named steps are Lean declarations),
\emph{Paper proof}, or \emph{External input}. Kernel-verified statements also receive written
proofs or proof outlines here.

\section{Discrepancy on the circle and orbits of rationals}\label{sec:prelim}

For a finite sample $u=(u_0,\ldots,u_{N-1})$ in $\T$, identified with $[0,1)$, let $\mu_u$ be its
empirical probability measure and $D(u)=\sup_{0\le a<c\le1}|\mu_u([a,c))-(c-a)|$. A
\emph{half-open arc} is the image in $\T$ of $[\alpha,\alpha+l)$ with $0\le l\le1$; its length is
$|I|=l$, and its \emph{excess} is $\mu_u(I)-|I|$.

\begin{lemma}[Arcs]\label{lem:arcs}\evidence{Paper proof}
For every half-open arc $I$, $|\mu_u(I)-|I||\le D(u)$, and $D(u)=\sup_I(\mu_u(I)-|I|)$, the supremum
of the excess over half-open arcs. Consequently $D(u)$ is unchanged if every $u_n$ is rotated by
the same $\theta\in\T$.
\end{lemma}

\begin{proof}
A non-wrapping arc is an interval $[a,c)$. A wrapping arc is the complement of an interval
$[c,a)$, and its excess is the negative of the excess of that interval; the full circle has
excess zero. This proves the first claim and that the supremum of excesses is at most $D(u)$.
Conversely an interval with negative excess $-e$ has a complementary arc with excess $e$.
Rotation permutes half-open arcs and preserves length.
\end{proof}

\begin{lemma}[Finite comparison]\label{lem:compare}\evidence{Paper proof}
Let $u,v$ be samples of length $N$, let $\epsilon\ge0$ and $E=\#\{n<N:\dist{u_n-v_n}>\epsilon\}$.
Then $|D(u)-D(v)|\le E/N+2\epsilon$.
\end{lemma}

\begin{proof}
Let $I=[\alpha,\alpha+l)$. If $l>1-2\epsilon$, then $\mu_u(I)-l<2\epsilon$. Otherwise each $n$
with $u_n\in I$ and $\dist{u_n-v_n}\le\epsilon$ has $v_n$ in the half-open arc
$I'=[\alpha-\epsilon,\alpha+l+\epsilon)$ of length $l+2\epsilon$. Hence
$\mu_u(I)\le\mu_v(I')+E/N\le l+2\epsilon+D(v)+E/N$ by \Cref{lem:arcs}. Taking the supremum
over $I$ gives $D(u)\le D(v)+E/N+2\epsilon$; exchange $u$ and $v$.
\end{proof}

When $v_n=\fr{b^n(x+z)}$ and $u_n=\fr{b^nx}$, the exceptional set depends only on $z$; this case
is the kernel-verified declaration
\lean{Phase18Synthetic}{583}{\ttt{discrepancy\_transfer}}. We shall often use the following
form, in which part of the orbit is left uncovered.

\begin{lemma}[Block certificate]\label{lem:blocks}\evidence{Paper proof}
Let $x\in\R$, $N\ge1$, and let $B_1,\ldots,B_m\subset[0,N)$ be disjoint sets of integers with
$U=N-\sum_i|B_i|$. Suppose that $y_n\in\T$ is given for $n\in\bigcup B_i$ and that for every $i$
and every half-open arc $I$,
$|\#\{n\in B_i:y_n\in I\}-|B_i||I||\le H_i$. If $\dist{b^nx-y_n}\le\epsilon$ for all covered $n$
except $E$ of them, then
\[
D_N(x)\le\frac{U+E+\sum_iH_i}{N}+2\epsilon .
\]
\end{lemma}

\begin{proof}
Put $y_n=\fr{b^nx}$ at uncovered $n$. For every half-open arc the count of $(y_n)_{n<N}$ differs
from $N|I|$ by at most $\sum_iH_i+U$, so $D(y)\le(U+\sum_iH_i)/N$ by \Cref{lem:arcs}. Apply
\Cref{lem:compare}.
\end{proof}

The rational models will be orbits of rationals with denominator prime to $b$. We use the
standard lifting of exponents: if $\ell$ is an odd prime and $\ell\mid y-1$, then
$v_\ell(y^m-1)=v_\ell(y-1)+v_\ell(m)$ for all $m\ge1$. (For $\ell\nmid m$ the quotient
$(y^m-1)/(y-1)\equiv m\pmod\ell$; and if $y=1+\ell^ju$ with $j\ge1$, then
$y^\ell=1+\ell^{j+1}u+\binom{\ell}{2}\ell^{2j}u^2+\cdots\equiv1+\ell^{j+1}u\pmod{\ell^{j+2}}$.)

\begin{lemma}[Grids]\label{lem:grid}\evidence{Lean-supported paper proof}
Let $p$ be an odd prime, $q=p^eQ$ with $\gcd(Q,pb)=1$ and $p\nmid b$. Let $\tau\ge1$ satisfy
$b^\tau\equiv1\pmod{pQ}$, and put $c=v_p(b^\tau-1)$. If $e\ge c$, then for every integer $a$ prime to
$p$ (not necessarily prime to $Q$) and every $j_0\ge0$, the $S=\tau p^{e-c}$ points $\fr{ab^j/q}$, $j_0\le j<j_0+S$, form $\tau$
arithmetic progressions in $\T$, each of $p^{e-c}$ points with common difference $p^{c-e}$.
Consequently, for every $L\ge1$, every $L$ consecutive points $\fr{ab^j/q}$, and every half-open arc
$I$,
\[
\bigl|\#\{j:\fr{ab^j/q}\in I\}-L|I|\bigr|\le Lp^{c-e}+\tau p^{e-c}.
\]
\end{lemma}

\begin{proof}
Put $y=b^\tau$. By lifting the exponent, $y^m\equiv1\pmod{p^e}$ if and only if
$p^{e-c}\mid m$, so $y$ generates a subgroup of order $p^{e-c}$ of the units modulo $p^e$, contained
in the set of residues congruent to $1$ modulo $p^c$; that set has $p^{e-c}$ elements, so the two
coincide. Since $a$ is a unit modulo $p^e$, the residues $ay^m$ modulo $p^e$ are exactly those
congruent to $a$ modulo $p^c$; since $y\equiv1\pmod Q$, $ay^m\equiv a\pmod Q$ for every integer $a$.
The Chinese remainder theorem therefore gives
$\{ay^m\bmod q\}=\{r\bmod q:r\equiv a\pmod{p^cQ}\}$, and these residues divided by $q$ form a
progression of $p^{e-c}$ points with difference $p^cQ/q$. Writing $j=j_0+r+\tau m$ with $0\le r<\tau$
splits $S$ consecutive points into $\tau$ such cosets. A progression of $M$ points with difference
$1/M$ meets every half-open arc $I$ in $M|I|+\theta$ points with $|\theta|\le1$. For general $L$
write $L=fS+g$ with $0\le g<S$; the error is at most $f\tau+g\le L\tau/S+S$. The order computation
and the count over one stretch of length $S$ starting at $j_0=0$ (equivalently, at any numerator)
are the Lean declarations \leanP{SurvivingGrid}{43}{\ttt{stride\_order}} and
\leanP{PrimePowerGrid}{337}{\ttt{crt\_period\_count}}, the latter also for numerators prime to $p$
only; the passage to arbitrary $L$ is written.
\end{proof}

The lemma counts along one chosen admissible stride $\tau$, which is all that \Cref{thm:Y} needs. For
an arbitrary admissible $\tau$ the $\tau$ progressions need not be distinct, and $p^{e-c}$ need not be
the largest size of an equally spaced subset of the orbit; we do not use the exact period or the
maximal grid.

\begin{corollary}[Prime-power orbits]\label{cor:period}\evidence{Paper proof}
Let $p\nmid b$ be an odd prime, $t=\ord_p(b)$ and $v=v_p(b^t-1)$. For $e\ge0$, $a$ prime to $p$,
$L\ge1$ and every half-open arc $I$, any $L$ consecutive points $\fr{ab^j/p^e}$ satisfy
\[
\bigl|\#\{j:\fr{ab^j/p^e}\in I\}-L|I|\bigr|\le p^vL/p^e+p^e .
\]
\end{corollary}

\begin{proof}
If $e\ge v$ apply \Cref{lem:grid} with $Q=1$, $\tau=t$, $c=v$, noting $tp^{e-v}<p^e$. If $e<v$ the
left side is at most $L\le p^vL/p^e$.
\end{proof}

The sharp form of \Cref{lem:compare} for translations is proved in the companion manuscript. We quote
it here and use it only in \Cref{prop:tails}.

\begin{theorem}[Sharp translation identity {\cite[Theorem~3.1]{CompanionH}}]\label{thm:sharpcited}\evidence{Kernel-verified}
For all real $x,z$ and every $N\ge1$, $|D_N(x+z)-D_N(x)|\le1-D_N(z)$, with equality at $x=0$.
\end{theorem}

This is Theorem~3.1 of \emph{The Arithmetic of Extreme Discrepancy: Translation, Addition Spectra and
Rational Operations} \cite{CompanionH}, in the same release, where it is proved; it is the Lean
declaration \lean{SharpTranslation}{343}{\ttt{sharp\_translation}}, with the equality case
\lean{SharpTranslation}{386}{\ttt{sharp\_translation\_isGreatest}}. We do not reprove it. Taking $x=-z$
gives $D_N(-z)=D_N(z)$.

\section{A worked example: the binary cubic number}\label{sec:example}

Let
\[
\xi=\sum_{k\ge1}\frac{1}{3^{v_3(k)}\,2^{2k^3}}
=\frac1{2^{2}}+\frac1{2^{16}}+\frac1{3\cdot2^{54}}+\frac1{2^{128}}+\cdots .
\]
Thus the coefficient $1/3^{v_3(k)}$ is inserted at binary position $2k^3$. The companion number
with all coefficients equal to one is $\xi_0=\sum_k2^{-2k^3}$, whose binary expansion consists of
isolated ones.

\begin{theorem}\label{thm:example}\evidence{Paper proof}
For every $N\ge1$, $D_N(\xi;2)\le200N^{-1/3}$. In particular $\xi$ is normal to base $2$. On the
other hand $D_N(\xi_0;2)\to1$.
\end{theorem}

\begin{proof}
Write $n_k=2k^3$ and, for $n_k\le n<n_{k+1}$, let
\[
R_n=2^n\sum_{j\le k}\frac{1}{3^{v_3(j)}2^{n_j}}
\]
be the rational number seen after $n$ shifts. Let $J=\lfloor\log_3k\rfloor$.

\emph{Step 1: the reduced denominator of $R_n$ is exactly $3^J$.} Every $j\le k$ has
$v_3(j)\le J$, so $3^JR_n=\sum_{j\le k}3^{J-v_3(j)}2^{n-n_j}$ is an integer. Modulo $3$ only the
indices with $v_3(j)=J$ survive, that is $j=3^Jm$ with $3\nmid m$ and $m\le k/3^J<3$; there are one
or two of them. Since $2^2\equiv1\pmod3$ and $n_j$ is even, each contributes $2^{n-n_j}\equiv2^n$.
Hence $3^JR_n\equiv(1\text{ or }2)\cdot2^n\not\equiv0\pmod3$.

\emph{Step 2: orbits of the models.} On $[n_k,n_{k+1})$ we have $R_{n+1}=2R_n$, so the points
$y_n=\fr{R_n}$ are consecutive points $\fr{a2^j/3^J}$ with $3\nmid a$. Here $t=\ord_3(2)=2$ and
$v_3(2^2-1)=1$. By \Cref{cor:period} the $\ell$ points of any piece of this block meet every
half-open arc in $\ell|I|+\theta$ points with $|\theta|\le3\ell/3^J+3^J$.

\emph{Step 3: the real orbit.} Positivity and distinctness of the positions give
$0\le2^n\xi-R_n\le\sum_{m\ge n_{k+1}}2^{n-m}=2^{1-(n_{k+1}-n)}$.

\emph{Step 4: counting.} Let $N\ge2$ and choose $K\ge1$ with $2K^3\le N<2(K+1)^3$. Apply
\Cref{lem:blocks} with blocks $B_k=[n_k,\min(n_{k+1},N))$, $1\le k\le K$, and $U=n_1=2$. The
block $B_k$ has length at most $h_k=n_{k+1}-n_k=2(3k^2+3k+1)$, and $k/3<3^{J}\le k$, so Step~2 gives
$H_k\le9h_k/k+k\le127k$ and $\sum_{k\le K}H_k\le127K^2$. Take
$W=\lceil\log_2(K+1)\rceil+2$ and $\epsilon=1/(2(K+1))$. By Step~3 the positions with
$n_{k+1}-n\ge W$ have $\dist{2^n\xi-y_n}\le2^{1-W}\le\epsilon$; at most $W$ positions in each
block fail this, so $E\le KW$. \Cref{lem:blocks} gives
\[
D_N(\xi)\le\frac{2+KW+127K^2}{N}+\frac1{K+1}.
\]
Using $N\ge2K^3$ and $W\le\log_2(K+1)+3\le4K$, the right side is at most
$(1+2+63.5+1)/K=67.5/K$. Since $N<2(K+1)^3\le16K^3$, $1/K<16^{1/3}N^{-1/3}<2.52N^{-1/3}$, and the
bound follows. For $N=1$ it is trivial.

For $\xi_0$ every $R_n$ is an integer and $0<\fr{2^n\xi_0}\le2^{1-(n_{k+1}-n)}$. Fix $W\ge1$. At
most $W(K+1)=O(N^{1/3})$ indices $n<N$ have $n_{k+1}-n<W$, so the interval $[2^{1-W},1)$ contains
$o(N)$ of the first $N$ points while having length $1-2^{1-W}$. Hence
$\liminf D_N(\xi_0)\ge1-2^{1-W}$ for every $W$.
\end{proof}

The constant $200$ is not optimized. For $N=2^{10},2^{11},\ldots,2^{18}$, enclosures of $D_N(\xi)$
to within $2^{-62}$ (computed exactly from the binary expansion and transferred by \Cref{lem:compare};
the script is listed in \Cref{sec:formal}) give values of $N^{1/3}D_N(\xi)$ between $3.477$, at
$N=2^{12}$, and $7.024$, at $N=2^{15}$. The proof uses only the denominator identity of Step~1 and complete periods. \Cref{sec:switch} shows
that Step~1 is an instance of a general identity, and that normality persists when the
complete-period count no longer gives a rate.

\section{Valuation-weighted expansions}\label{sec:switch}

Fix a prime $p\nmid b$, $t=\ord_p(b)$, an integer $s\ge0$, and strictly increasing positive integers
$(n_k)_{k\ge1}$ with $t\mid n_k$. Put $n_0=0$, $w(k)=p^{-sv_p(k)}$,
\[
X=X_s=\sum_{k\ge1}w(k)b^{-n_k},\qquad U_K=\sum_{k\le K}w(k)b^{-n_k},
\]
and for $n\ge0$ let $K(n)=\#\{k\ge1:n_k\le n\}$ and $R_n=b^nU_{K(n)}$.

\subsection{Grouped and raw forms}

A sparse expansion $\sum_kc_kb^{-n_k}$ also has a dense representation, in which the coefficient
$c_k$ is spread over the whole gap $(n_{k-1},n_k]$. When the multipliers have numerator one and the
gaps are divisible by every positive integer, the reduced denominators of the dense coefficients
involve infinitely many primes, which disappear in the sparse form.

\begin{proposition}[Geometric grouping]\label{prop:group}\evidence{Kernel-verified}
For integers $m\ge0$, $h\ge1$ and every field element $b$ with $b\ne0$ and $b^h\ne1$,
\[
\frac{b-1}{b^h-1}\sum_{r=1}^{h}b^{-(m+r)}=b^{-(m+h)} .
\]
Consequently, for real $b>1$, every strictly increasing $(n_k)$ with $n_0=0$ and bounded real
multipliers $c_k$, the series whose coefficient at $b^{-m}$, $n_{k-1}<m\le n_k$, equals
$c_k(b-1)/(b^{n_k-n_{k-1}}-1)$ converges absolutely to $\sum_kc_kb^{-n_k}$. Write
$\rho_h=(b^h-1)/(b-1)$ for the repunit of length $h$. If $b\ge2$ is an integer and for every $M\ge1$
some gap $n_k-n_{k-1}$ is divisible by $M$ (for example $n_k=t\,k!$), then every prime not dividing $b$
divides $\rho_{n_k-n_{k-1}}$ for some $k$, and in particular infinitely many primes divide the gap
repunits $\rho_{n_k-n_{k-1}}$.
\end{proposition}

\begin{proof}
The first identity is the finite geometric sum
$\sum_{r=1}^hb^{-(m+r)}=b^{-(m+h)}(b^h-1)/(b-1)$. Absolute convergence follows because every raw
coefficient is bounded by $\sup|c_k|$, and grouping an absolutely convergent series preserves its
sum. Let $r\nmid b$ be prime. If $r\nmid b-1$, then $r\mid\rho_{r-1}$ by Fermat's theorem; if
$r\mid b-1$, then $\rho_r\equiv r\equiv0\pmod r$. In either case $r\mid\rho_{r(r-1)}$, and
$\rho_a\mid\rho_h$ whenever $a\mid h$. Lean declarations:
\lean{Phase18Synthetic}{76}{\ttt{geometric\_retirement\_weight}},
\lean{Phase18Synthetic}{146}{\ttt{grouped\_retirement\_hasSum}},
\lean{Phase18Synthetic}{290}{\ttt{raw\_prime\_support\_infinite}}.
\end{proof}

The Lean statement concerns the repunits, not the reduced raw coefficients, and the passage between
the two needs the form of the multipliers. If $c_k=1/m_k$ with a positive integer $m_k$, as for the
weights $p^{-sv_p(k)}$, then $b-1\mid b^h-1$ gives the raw coefficient $1/(m_k\rho_h)$, $h=n_k-n_{k-1}$,
whose reduced denominator $m_k\rho_h$ is divisible by $\rho_h$; so under the gap hypothesis every prime
not dividing $b$ occurs in some reduced raw denominator (a written step, not formalized). For general
multipliers this fails: with $c_k=\rho_h b^{-h}\in(0,1]$ every raw coefficient on the gap equals
$b^{-h}$, and only primes dividing $b$ occur. The gap hypothesis is also needed: a schedule with
finitely many distinct gaps involves only finitely many repunits.

\subsection{Exact denominators and the switch}

\begin{theorem}[Exact reduced denominators]\label{thm:den}\evidence{Kernel-verified}
For every $K\ge0$ and every $n\ge n_K$, the reduced denominator of $b^nU_K$ is
$p^{s\lfloor\log_pK\rfloor}$, where $\lfloor\log_p0\rfloor=0$. In particular the reduced denominator
of $R_n$ is $p^{s\lfloor\log_pK(n)\rfloor}$ for every $n\ge0$.
\end{theorem}

\begin{proof}
Let $J=\lfloor\log_pK\rfloor$ and $K\ge1$, $s\ge1$ (otherwise the value is an integer). Every
$k\le K$ has $v_p(k)\le J$, so $p^{sJ}b^nU_K=\sum_{k\le K}p^{s(J-v_p(k))}b^{n-n_k}$ is an integer.
Modulo $p$ only the indices $k=p^Jm$ with $p\nmid m$ survive, and $m\le K/p^J<p$, so there are
between $1$ and $p-1$ of them. Since $t\mid n_k$, each contributes $b^{n-n_k}\equiv b^n\not\equiv0$.
The count is a unit modulo $p$, so the numerator is prime to $p$. Lean declarations:
\lean{RetirementDenominator}{110}{\ttt{shiftedPrefix\_den}},
\lean{RetirementDenominator}{157}{\ttt{fullPrefix\_den}},
\lean{RetirementDenominator}{67}{\ttt{maximalLayer\_card}}.
\end{proof}

The qualitative analytic input is a fixed-support estimate for short orbits.

\begin{lemma}[Short orbits]\label{lem:short}\evidence{Kernel-verified}
Let $b,c\ge2$ be coprime integers. For every integer $d_0\ge0$ and every $\eta>0$ there is $M$ such
that for every rational $y$ whose reduced denominator $q$ divides a power of $c$ and satisfies
$M\le q\le T^{d_0}$,
\[
\Bigl|\sum_{j=0}^{T-1}e(b^jy)\Bigr|\le\eta T,\qquad e(\theta)=\exp(2\pi i\theta).
\]
\end{lemma}

This is \lean{SmoothOrbit}{207}{\ttt{fixed\_support\_orbit\_estimate}} (definition at
\lean{SmoothOrbit}{16}{\ttt{SmoothOrbitEstimate}}), proved there by elementary differencing. It is
also a qualitative consequence of Vandehey's quantitative bound \cite[Corollary~1.2]{Vandehey19}:
for a finite prime set $P$, $m$ supported on $P$, $\gcd(m,b)=\gcd(a,m)=1$, $\epsilon>0$ and
$T\ge m^\epsilon$, one has $|\sum_{j=1}^{T}e(ab^j/m)|\le C\,m^{-\delta}T(1+\log m)$ with
$\delta=\delta(\epsilon)>0$ and $C=C(\epsilon,b,P)$.


The short-orbit estimate gives normality from denominator information alone, at every prefix.

\begin{proposition}[All-prefix criterion]\label{prop:criterion}\evidence{Lean-supported paper proof}
Let $x=\sum_{m\ge0}a_mb^{-m}$ with rational $a_m$ satisfying $|a_m|\le C_0(m+1)^A$ for constants
$C_0,A\ge0$, put $R_n=\sum_{m\le n}a_mb^{n-m}$, and let $q_n$ be the reduced denominator of $R_n$.
Suppose that for some $\eta\in(0,1]$, some integer $c\ge2$ prime to $b$ and some $D_0>0$:
\begin{enumerate}[label=\textup{(\roman*)},nosep]
\item $\#\{m\le N:a_m\ne0\}=O(N^{1-\eta})$;
\item every $q_n$ divides a power of $c$, and $q_n\le C_0(n+1)^{D_0}$;
\item for every fixed $M$, $\#\{n<N:q_n<M\}=o(N)$.
\end{enumerate}
Then $x$ is normal to base $b$.
\end{proposition}

\begin{proof}
The Lean-supported step is \Cref{lem:short}. Fix an integer $h\ne0$ and $\eta_0>0$. For large $N$ put
$T=\lfloor N^{\eta/2}\rfloor$ and $W=\lceil(A+3)\log_b(2N)\rceil\le T$, and split $[0,N)$ into consecutive
blocks $[a,a+T)$, the last possibly shorter. Call a block \emph{clean} if no $m$ with $a_m\ne0$ lies in
$(a,a+T+W]$. Each such $m$ makes at most two blocks unclean, so by (i) the unclean blocks and the last
block cover $O(TN^{1-\eta})+T=o(N)$ indices.

On a clean block and for $0\le j<T$ we have $R_{a+j}=b^jR_a$, and the omitted terms of $b^{a+j}x$ have
positions $m>a+j+W$, so
\[
|b^{a+j}x-R_{a+j}|\le C_0\sum_{i>W}(2N+i)^Ab^{-i}\ll_{A,b,C_0}N^AW^Ab^{-W}=o(1).
\]
Since $\gcd(b,c)=1$, all $q_{a+j}$ equal $q_a$, and the reduced denominator $q'$ of $hR_a$ divides a power
of $c$ and satisfies $q_a/|h|\le q'\le q_a$. Choose an integer $d_0>2D_0/\eta+1$, so that
$q_a\le C_0(N+1)^{D_0}\le T^{d_0}$ for large $N$, and let $M_0$ be given by \Cref{lem:short} for $b$, $c$,
$d_0$ and $\eta_0$. On clean blocks with $q'\ge M_0$, $|\sum_{j<T}e(b^jhR_a)|\le\eta_0T$. Clean blocks with
$q'<M_0$ have $q_n<|h|M_0$ at each of their indices, so by (iii) they cover $o(N)$ indices. Hence
\[
\Bigl|\sum_{n<N}e(hb^nx)\Bigr|\le o(N)+\eta_0N+2\pi|h|N\cdot o(1),
\]
and $\limsup_N|N^{-1}\sum_{n<N}e(hb^nx)|\le\eta_0$ for every $\eta_0>0$. Weyl's criterion
\cite[Chapter~1, Theorem~2.1, p.~7]{KN} and \cite[Chapter~1, Theorem~8.1, p.~70]{KN} give normality.
\end{proof}

\begin{theorem}[Normality switch]\label{thm:switch}\evidence{Kernel-verified}
Assume in addition that $n_k\ge k^2$ for all $k$. Then $D_N(X_0)\to1$, and for every $s\ge1$,
$D_N(X_s)\to0$; equivalently $X_s$ is normal to base $b$ if and only if $s\ge1$. In particular this
holds for $n_k=tk^d$ with any integer $d\ge2$.
\end{theorem}

\begin{proof}
The statement is a Lean theorem (declarations listed at the end of this proof), with the definitions
of discrepancy and normality of \Cref{sec:intro}. We outline the argument.

Since $n_k\ge k^2$, $K(M)\le\sqrt M$. Positivity gives
$0\le b^nX-R_n\le\frac{b}{b-1}b^{-(n_{K(n)+1}-n)}$, and for fixed $W$ at most
$W(K(N+W)+1)=o(N)$ indices $n<N$ have an insertion in $(n,n+W]$. If $s=0$, every $R_n$ is an integer,
so all but $o(N)$ of the points $\fr{b^nX_0}$ lie in $[0,\frac{b}{b-1}b^{-W})$; as in the proof of
\Cref{thm:example}, $\liminf D_N(X_0)\ge1-\frac b{b-1}b^{-W}$ for every $W$. If $s\ge1$, apply
\Cref{prop:criterion} with $c=p$ and $\eta=1/2$: the coefficients are at most one, there are at most
$\sqrt N$ insertions below $N$, and by \Cref{thm:den} the reduced denominator of $R_n$ is
$p^{s\lfloor\log_pK(n)\rfloor}\le(n+1)^s$, which tends to infinity with $n$. For $n_k=tk^d$ the hypotheses
hold because $t\ge1$ and $d\ge2$. Lean declarations:
\lean{RetirementDenominator}{436}{\ttt{discrepancy\_switch}},
\lean{RetirementDenominator}{444}{\ttt{normal\_iff\_pos}} and, for $n_k=tk^d$,
\lean{RetirementDenominator}{456}{\ttt{polynomial\_discrepancy\_switch}}.
\end{proof}

\begin{corollary}[A normal Liouville family]\label{cor:liouville}\evidence{Lean-supported paper proof}
For every $s\ge1$, the number $L=\sum_{k\ge1}p^{-sv_p(k)}b^{-t(k+1)!}$ is normal to base $b$ and is a
Liouville number; in particular it is transcendental.
\end{corollary}

\begin{proof}
The schedule $n_k=t(k+1)!$ consists of strictly increasing multiples of $t$ with $n_k\ge k^2$, so
normality is the kernel-verified \Cref{thm:switch} (the verification of these two elementary
hypotheses is not formalized). The partial sum $U_k$ has reduced denominator
$Q_k=b^{n_k}p^{s\lfloor\log_pk\rfloor}\le b^{n_k}k^s$: indeed the newest term is the only one not
divisible by $b$ after clearing $b^{n_k}$, and \Cref{thm:den} gives the $p$-part. Moreover
$0<L-U_k\le\frac{b}{b-1}b^{-n_{k+1}}=\frac b{b-1}b^{-(k+2)n_k}$. For each fixed $\mu>0$ this is
smaller than $Q_k^{-\mu}$ for all large $k$, and the $U_k$ are distinct. Hence $L$ is a Liouville
number, and transcendental by Liouville's theorem.
\end{proof}

\subsection{The complete-period rate}

\begin{theorem}[Complete-period rate]\label{thm:rate}\evidence{Lean-supported paper proof}
Let $p\nmid b$ be an odd prime, $v=v_p(b^t-1)$, $s\ge1$, $d\ge s+2$, and $n_k=tk^d$. There is
$C=C(b,p,s,d)$ such that for all $N\ge2$,
\[
D_N(X_s)\le C\bigl(N^{-s/d}+N^{-(d-s-1)/d}+N^{-(d-1)/d}\log N\bigr).
\]
The exponent $\min(s,d-1-s)/d$ equals $1/3$ for $(s,d)=(1,3)$ and tends to $1/2$ along
$(s,d)=(s,2s+1)$.
\end{theorem}

\begin{proof}
The Lean-supported step is the denominator identity, \Cref{thm:den}; the rest follows the proof of
\Cref{thm:example}. Let $tK^d\le N<t(K+1)^d$ with $K\ge1$ (smaller $N$ are absorbed in $C$). On the
block $B_k=[n_k,\min(n_{k+1},N))$ the model $y_n=\fr{R_n}$ runs through consecutive points of
the orbit of a rational with reduced denominator $q_k=p^{s\lfloor\log_pk\rfloor}\in((k/p)^s,k^s]$.
The block has length at most $h_k=n_{k+1}-n_k\le t\,d\,2^{d-1}k^{d-1}$. By
\Cref{cor:period},
\[
\sum_{k\le K}H_k\le\sum_{k\le K}\Bigl(\frac{p^vh_k}{q_k}+q_k\Bigr)
\ll\sum_{k\le K}\bigl(k^{d-1-s}+k^s\bigr)\ll K^{d-s}+K^{s+1},
\]
using $d-1-s\ge1$. Take $W=\lceil2\log_b(2N)\rceil$ and $\epsilon=\frac b{b-1}b^{-W}\le2N^{-2}$; the
exceptional count is $E\le KW$, and $U=n_1=t$. \Cref{lem:blocks} and $N\asymp K^d$ give
$D_N\ll K^{-s}+K^{-(d-s-1)}+K^{1-d}\log N+N^{-2}$, which is the claim.
\end{proof}

\begin{remark}
Two refinements are not claimed here. First, bounding the final incomplete period on each block by
Gauss sums instead of trivially improves the range of $d$ for which a power saving follows from
complete periods. Second, Vandehey's \cite[Corollary~1.2]{Vandehey19}, applied on blocks of length
$N^\beta$ with $0<\beta<1-1/d$, should give some power saving for every $s\ge1$ and $d\ge2$. We have
not written out either argument in full, and in particular we have not checked the dependence of
Vandehey's constants on the block parameters; we therefore state no rate outside
$d\ge s+2$.
\end{remark}

\section{Prescribing the discrepancy limit}\label{sec:program}

\begin{lemma}[Positive state programming]\label{lem:program}\evidence{Kernel-verified}
Let $0=n_0<n_1<\cdots$ be integers and $u_0=0,u_1,u_2,\ldots$ rationals in $[0,1)$. Define $R_0=0$,
$c_k\in(0,1]$ as the representative of $u_k-b^{n_k-n_{k-1}}R_{k-1}$ modulo $1$, and
$R_k=b^{n_k-n_{k-1}}R_{k-1}+c_k$. Then $R_k=b^{n_k}\sum_{i\le k}c_ib^{-n_i}$ and $\fr{R_k}=u_k$ for
every $k$. If every $u_k$ has reduced denominator a power of a prime $p$, so has every $c_k$.
\end{lemma}

\begin{proof}
By induction $R_k-u_k=b^{n_k-n_{k-1}}(R_{k-1}-u_{k-1})+(b^{n_k-n_{k-1}}u_{k-1}+c_k-u_k)\in\Z$. The
coefficient $c_k$ differs by an integer from $u_k-b^{n_k-n_{k-1}}u_{k-1}$, whose denominator divides
$\den(u_k)\den(u_{k-1})$. Lean declarations:
\lean{Phase18Synthetic}{383}{\ttt{programState\_fract}},
\lean{Phase18Synthetic}{438}{\ttt{programState\_eq\_prefix}},
\lean{Phase18Synthetic}{459}{\ttt{programmed\_prefix\_fract}},
\lean{Phase18Synthetic}{364}{\ttt{programCoeff\_mem}},
\lean{Phase18Synthetic}{426}{\ttt{programCoeffs\_prime\_power}}.
\end{proof}

For $x=\sum_kc_kb^{-n_k}$ as in \Cref{lem:program} and $n_k\le n<n_{k+1}$, the rational part after
$n$ shifts is $b^{n-n_k}R_k$, whose fractional part is $y_n=\fr{b^{n-n_k}u_k}$. Since
$0<c_i\le1$ and the positions are distinct,
\begin{equation}\label{eq:tail}
0<b^nx-b^{n-n_k}R_k\le\frac b{b-1}b^{-(n_{k+1}-n)} .
\end{equation}
Call block $k$ (the integers in $[n_k,n_{k+1})$) \emph{retired} if $u_k=0$ and \emph{active}
otherwise; block $0$ is retired.

\begin{proposition}[Modelled discrepancy]\label{prop:model}\evidence{Paper proof}
Let $n_K\le N<n_{K+1}$, let $Z$ be the number of indices $n<N$ in retired blocks, and suppose that
for every half-open arc $I$ the active indices $n<N$ satisfy
$|\#\{n:y_n\in I\}-(N-Z)|I||\le H$. Then for every integer $W\ge0$,
\[
\Bigl|D_N(x)-\frac ZN\Bigr|\le\frac{H+(K+1)W}{N}+\frac{2b}{b-1}b^{-W}.
\]
\end{proposition}

\begin{proof}
Retired models are $0$. For every half-open arc,
$\#\{n<N:y_n\in I\}-N|I|=(\text{active error})+Z(\mathbf 1_{0\in I}-|I|)$, which has modulus at most
$H+Z$, and equals at least $Z(1-|I|)-H$ for $I=[0,|I|)$. Letting $|I|\to0$ and using
\Cref{lem:arcs}, $|D(y)-Z/N|\le H/N$. By \eqref{eq:tail} at most $W$ indices of each block
$k\le K$ have $\dist{b^nx-y_n}>\frac b{b-1}b^{-W}$, and \Cref{lem:compare} completes the proof.
\end{proof}

\begin{theorem}[Every discrepancy limit]\label{thm:program}\evidence{Lean-supported paper proof}
Let $p\nmid b$ be an odd prime, $t=\ord_p(b)$ and $v=v_p(b^t-1)$. For $\delta\in[0,1]$ put $n_k=tk^3$,
$a_k=\lfloor k\delta\rfloor-\lfloor(k-1)\delta\rfloor\in\{0,1\}$, $q_k=p^{1+\lfloor\log_pk\rfloor}$, and
$u_k=0$ if $a_k=1$, $u_k=1/q_k$ if $a_k=0$. Let $x_\delta=\sum c_kb^{-n_k}$ be given by
\Cref{lem:program}. Then every $c_k\in(0,1]$ has reduced denominator dividing $q_k$, so a power of $p$
at most $pk$, and for all $N\ge1$,
\[
|D_N(x_\delta)-\delta|\le C_{b,p}N^{-1/3}
\]
with $C_{b,p}$ independent of $\delta$. Moreover the orbit measures
$N^{-1}\sum_{n<N}\delta_{\fr{b^nx_\delta}}$ converge weakly to $(1-\delta)\lambda+\delta\,\delta_0$, where
$\lambda$ is Lebesgue measure, and $x_\delta$ is irrational.
\end{theorem}

\begin{proof}
The Lean-supported step is \Cref{lem:program}. Since $q_{k-1}\mid q_k$, the coefficient $c_k$, which
differs by an integer from $u_k-b^{n_k-n_{k-1}}u_{k-1}$, has denominator dividing $q_k\le pk$. Let
$N\ge t$ and $tK^3\le N<t(K+1)^3$. On an active
block $k$ the models are consecutive points $\fr{b^j/q_k}$, with $k<q_k\le pk$, and
$h_k=n_{k+1}-n_k=t(3k^2+3k+1)\le7tk^2$. \Cref{cor:period} gives
$H\le\sum_{k\le K}(p^vh_k/q_k+q_k)\le(7tp^v+p)K(K+1)/2$, independently of the mask $(a_k)$.

The retired count is $Z=t+\sum_{1\le k<K}a_kh_k+a_K(N-n_K)$, and $N=t+\sum_{1\le k<K}h_k+(N-n_K)$. Put
$S_j=\sum_{i\le j}(a_i-\delta)$, so $|S_j|\le1$. Summation by parts with the increasing weights $h_k$
gives $|\sum_{1\le k<K}(a_k-\delta)h_k|\le2h_{K-1}$, and therefore
\begin{equation}\label{eq:Z}
|Z-\delta N|\le t+3h_K .
\end{equation}
\Cref{prop:model} with $W=\lceil\log_b(\frac b{b-1}(K+1))\rceil$ yields
\[
|D_N(x_\delta)-\delta|\le\frac{(7tp^v+p)K(K+1)/2+(K+1)W+t+3h_K}{N}+\frac2{K+1},
\]
which is $O_{b,p}(K^{-1})=O_{b,p}(N^{-1/3})$, uniformly in $\delta$. For $N<t$ use $D_N\le1$.

For weak convergence let $f$ be continuous on $\T$ and approximate it uniformly by step functions
on a fixed finite partition into arcs. The active indices have counts $(N-Z)|I|+O(H)$ on each arc,
the retired real points lie within $\frac b{b-1}b^{-W}$ of $0$ off $(K+1)W$ indices, and
$Z/N\to\delta$ by \eqref{eq:Z}. If $\delta<1$ the limit has a nonzero absolutely continuous part,
whereas the orbit measures of a rational number converge to the uniform measure on a finite cycle;
so $x_\delta$ is irrational. If $\delta=1$ every $u_k$ is $0$, every $c_k=1$, and the base-$b$ expansion
of $x_1$ has isolated digits $1$ separated by unbounded gaps, which is not eventually periodic.
\end{proof}

For example, with $b=2$, $p=3$ and $\delta=2/5$, the first five insertions are at
$2,16,54,128,250$, the targets are $1/3,1/3,0,1/9,0$, and the coefficients are
$1/3,1,2/3,1/9,5/9$.

\begin{remark}
The canonical numbers $x_\delta$ are irrational, but we do not know their algebraic nature. Altering
the mask at the sparse indices $k=2^j$ does not change the rate and produces continuum many
distinct realizations of each $\delta$; all but countably many are transcendental.
\end{remark}

\section{Lower and upper discrepancies and invariant measures}\label{sec:limitsets}

\begin{theorem}[All pairs of extreme limits]\label{thm:pairs}\evidence{Lean-supported paper proof}
Let $p,t,n_k,q_k$ be as in \Cref{thm:program} and $0\le\alpha\le\beta\le1$. There are rationals
$c_k\in(0,1]$ with reduced denominators dividing $q_k\le pk$, given by \Cref{lem:program} from targets
$u_k\in\{0,1/q_k\}$ and an explicit mask, such that $x=\sum_kc_kb^{-n_k}$ satisfies
$\liminf_ND_N(x)=\alpha$ and $\limsup_ND_N(x)=\beta$. More precisely $D_N(x)=\Lambda_N/N+O(N^{-1/3})$,
with $\Lambda_N$ defined in the proof.
\end{theorem}

\begin{proof}
Let $K_1=2$, $K_{j+1}=K_j^2$, and for $K_j\le k<K_{j+1}$ let $\theta_k=\alpha$ for odd $j$ and
$\theta_k=\beta$ for even $j$; put $\theta_1=\alpha$. Use the mask
$a_k=\lfloor\Theta_k\rfloor-\lfloor\Theta_{k-1}\rfloor$ with $\Theta_k=\sum_{i\le k}\theta_i$, and the
targets of \Cref{thm:program}. Then $|\sum_{i\le j}(a_i-\theta_i)|\le1$, and the argument for
\eqref{eq:Z} gives $|Z-\Lambda_N|\le t+3h_K$, where $\Lambda_N=\sum_{k\le K}\theta_k\ell_k$ and $\ell_k$
is the length of the part of block $k$ below $N$. Since the active-block bound is independent of
the mask, \Cref{prop:model} gives $D_N(x)=\Lambda_N/N+O(N^{-1/3})$. Each $\Lambda_N/N$ is a weighted
average of the $\theta_k$ up to $O(t/N)$, so lies in $[\alpha,\beta]+o(1)$. At $N=tK_{j+1}^3$ the
blocks before stage $j$ have total length at most $tK_j^3=o(N)$, so $\Lambda_N/N$ tends to $\alpha$ along
odd $j$ and to $\beta$ along even $j$.
\end{proof}

Let $\Prob$ be the space of Borel probability measures on $\T$ with the weak topology and $\Minv$
the closed subset of $T$-invariant measures. Fix nonzero trigonometric polynomials $f_1,f_2,\ldots$
with rational coefficients, dense in $C(\T)$, and put
$d(\mu,\nu)=\sum_i2^{-i}|\int f_id\mu-\int f_id\nu|/\|f_i\|_\infty$. Uniform approximation shows that
$d(\mu_n,\mu)\to0$ if and only if $\mu_n\to\mu$ weakly, and $\Prob$ is sequentially compact by a
diagonal argument and the Riesz representation theorem. Since each summand of $d$ is a convex
function of $\mu$, $d(\sum_r\gamma_r\mu_r,\nu)\le\max_rd(\mu_r,\nu)$ for convex weights $\gamma_r$. For $x\in\R$ let
$\mu_N^x=N^{-1}\sum_{n<N}\delta_{\fr{b^nx}}$ and $V(x)$ its set of weak limit points.

\begin{lemma}\label{lem:necessity}\evidence{Paper proof}
For every $x$, $V(x)$ is a nonempty compact connected subset of $\Minv$.
\end{lemma}

\begin{proof}
Compactness of $\Prob$ gives nonemptiness, and $V(x)$ is closed. For continuous $f$,
$\int f\circ T\,d\mu_N^x-\int fd\mu_N^x=(f(T^Nx)-f(x))/N\to0$, so limits are invariant. Since
$d(\mu_N^x,\mu_{N+1}^x)\le2/(N+1)$, connectedness follows from the general fact that the cluster set
of a sequence $z_N$ in a compact metric space with $d(z_N,z_{N+1})\to0$ is connected: if it were
the union of disjoint nonempty closed sets $A,B$ at distance $3r>0$, the sequence would eventually
avoid the set of points at distance at least $r$ from $A\cup B$ (a compact set containing no
cluster point), yet would pass infinitely often from the $r$-neighbourhood of $A$ to that of $B$ in
steps smaller than $r$, which is impossible.
\end{proof}

For $\mu\in\Prob$ let $D(\mu)=\sup_I|\mu(I)-|I||$ over half-open arcs; for an empirical measure this
is the discrepancy of \Cref{sec:prelim}.

\begin{lemma}\label{lem:Dcont}\evidence{Paper proof}
If $\mu_i\to\mu$ weakly in $\Prob$, then $D(\mu_i)\to D(\mu)$, even when $\mu$ has atoms.
\end{lemma}

\begin{proof}
As in \Cref{lem:arcs}, $D(\nu)$ is the supremum of the excess $\nu(I)-|I|$ over half-open arcs. Given
$\eta>0$, choose a rotated grid of $M\ge1/\eta$ equally spaced points that avoids the countably many
atoms of $\mu$. Every half-open arc is contained in an arc whose endpoints are grid points (or in the
whole circle), at a cost of at most $2\eta$ in length, so for every $\nu$ the maximum excess over the
finitely many grid arcs lies in $[D(\nu)-2\eta,D(\nu)]$. The masses of grid arcs converge, because
their endpoints are not atoms of $\mu$. Hence $\limsup|D(\mu_i)-D(\mu)|\le4\eta$.
\end{proof}

\begin{theorem}[Invariant-measure limit sets]\label{thm:limitsets}\evidence{Lean-supported paper proof}
A set $F\subset\Prob$ equals $V(x)$ for some real $x$ if and only if $F$ is a nonempty compact
connected subset of $\Minv$. For every integer $t\ge1$, every such $F$ is $V(x)$ for some
$x=\sum_{k\ge1}c_kb^{-tk^3}$ with rational $c_k\in(0,1]$ given by \Cref{lem:program}, whose state $u_k$
after the $k$th insertion is $0$ or a purely periodic rational with period at most $k$.
\end{theorem}

\begin{proof}
Necessity is \Cref{lem:necessity}. The Lean-supported step of the sufficiency proof is
\Cref{lem:program}, which programs the states.

\emph{Periodic approximation.} Let $\mu\in\Minv$, $\eta>0$, and let $g_1,\ldots,g_r$ be
Lipschitz functions with constant at most $\Lambda$ and $\|g_i\|_\infty\le1$. We find a rational $y$
with purely periodic base-$b$ expansion whose cycle measure $\pi$ satisfies
$|\int g_id\pi-\int g_id\mu|<\eta$. Fix integers $m$ and $L$ with $\Lambda b^{-m}+2m/L<\eta/3$, and let
$A_Lg(z)=L^{-1}\sum_{j<L}g(T^jz)$. By invariance $\int g_id\mu=\int A_Lg_i\,d\mu$. Partition the
range of $(A_Lg_1,\ldots,A_Lg_r)$ into Borel boxes of diameter $<\eta/3$, choose a point $z_\iota$ in
each preimage of positive measure, and replace the weights by rationals $M_\iota/M$ summing to one,
so that $|\sum_\iota(M_\iota/M)A_Lg_i(z_\iota)-\int g_id\mu|<2\eta/3$. Let $y$ be the purely periodic
number whose period is the concatenation, over $\iota$, of $M_\iota$ copies of the first $L$
base-$b$ digits of $z_\iota$. If a copy of the word of $z_\iota$ starts at offset $o$ of the period
and $0\le j\le L-m$, then $T^{o+j}y$ and $T^jz_\iota$ share their first $m$ digits, so
$\dist{T^{o+j}y-T^jz_\iota}\le b^{-m}$. The remaining $m$ offsets of each copy change $g_i$ by at
most $2$. Averaging over the period gives the claim; an
all-$(b-1)$ period represents $0\in\T$ and the same estimate holds on $\T$.

\emph{Chains.} Since $F$ is compact and connected, for every $\eta>0$ any two points of $F$ are
joined by a finite chain in $F$ with steps less than $\eta$ (the set of points reachable from a
given point is open and closed in $F$). Concatenating, for $r=1,2,\ldots$, chains with steps
$<1/r$ that start and end at a fixed point and visit a $1/r$-net of $F$, we obtain $\mu_j\in F$
with $d(\mu_j,\mu_{j+1})\to0$ whose cluster set is $F$.

\emph{Construction.} By the first step, applied to $g_i=f_i/\|f_i\|_\infty$ for $i\le j$, choose
periodic rationals $y_j\in[0,1)$ with cycle measures $\pi_j$, periods $P_j$, and
$d(\pi_j,\mu_j)<1/j+2^{1-j}$. Choose $K_1\ge\max(2,P_1)$ and $K_{j+1}\ge\max(K_j^2,P_{j+1})$. Let
$n_k=tk^3$, $u_k=0$ for $k<K_1$ and $u_k=y_j$ for $K_j\le k<K_{j+1}$, and let $x$ be given by
\Cref{lem:program}; then $u_k$ has period $P_j\le K_j\le k$. On block $k$ in stage $j$ the models are
consecutive points of the cycle of $y_j$; for a Lipschitz $g$ with $\|g\|_\infty\le1$, a piece of length
$\ell_k$ satisfies $|\sum g(y_n)-\ell_k\int gd\pi_j|\le2P_j\le2k$. Summing over $k\le K$, where
$tK^3\le N<t(K+1)^3$, gives $O(K^2)=O_t(N^{2/3})$. By
\eqref{eq:tail}, the real points are within $\frac b{b-1}b^{-W}$ of the models off $(K+1)W$ indices.
Hence $d(\mu_N^x,\nu_N)\to0$, where $\nu_N=N^{-1}\sum_k\ell_k\pi_{j(k)}$ (with $\pi=\delta_0$ before
stage $1$).

\emph{Limits.} Let $N_j=tK_{j+1}^3$ be the end of stage $j$. Stages before $j$ occupy the indices
below $tK_j^3$ and so have weight at most $(K_j/K_{j+1})^3\le K_j^{-3}$ in $\nu_{N_j}$, whence
$d(\nu_{N_j},\mu_j)\to0$. For $N_{j-1}<N\le N_j$,
$\nu_N$ is a convex combination of $\nu_{N_{j-1}}$ and $\pi_j$, both within $o(1)$ of $\mu_j$ because
$d(\mu_{j-1},\mu_j)\to0$; by convexity of balls $d(\nu_N,\mu_j)\to0$. Therefore every limit point of
$\mu_N^x$ lies in $F$, and every point of $F$, being a limit of a subsequence of $(\mu_j)$, is a limit
of the corresponding $\mu_{N_j}^x$.
\end{proof}

The symbolic realization of every nonempty compact connected set of invariant measures goes back to
Colebrook \cite{Colebrook} for base-$b$ expansions and to Sigmund \cite[Theorem~4]{Sigmund74} for systems
with the specification property; see also \cite{DGS}. The density of periodic orbit measures used in
the first step goes back to Parthasarathy \cite{Parthasarathy61}. The next remark shows that the
restriction to positive rational coefficients on a fixed schedule is, by itself, no restriction.

\begin{remark}[Qualitative content]\label{rem:qualitative}
Let $p$ be a prime and $0<n_1<n_2<\cdots$ any integers. For every $y\in[0,1)$ there are $m\ge0$ and
rationals $c_k\in(0,1)$ whose denominators are powers of $p$ such that
$x=\sum_kc_kb^{-n_k}$ satisfies $b^mx\equiv y\pmod1$. Indeed choose $m$ with
$x=(1+y)b^{-m}<b^{-n_1}$. If $0<r<b^{-n_k}$ is the residual before step $k$ (initially $r=x$), the
open interval of $c$ with $0<r-cb^{-n_k}<b^{-n_{k+1}}$ and $0<c<b^{n_k}r$ is nonempty and contains a
rational with $p$-power denominator; take it as $c_k$. Then $c_k<1$, the new residual lies in
$(0,b^{-n_{k+1}})$, and the residuals tend to $0$, so $x=\sum_kc_kb^{-n_k}$. The orbit of $x$ agrees
with that of $y$ after $m$ shifts, so $D_N(x)-D_N(y)=O(m/N)$ and $V(x)=V(y)$. Consequently, by the
classical realization above and \Cref{lem:Dcont}, every nonempty compact connected $F\subset\Minv$ is
$V(x)$ for such an $x$ on any schedule, even with a single denominator prime; taking
$F=\{(1-\theta)\lambda+\theta\delta_0:\alpha\le\theta\le\beta\}$, on which $D$ takes the values
$[\alpha,\beta]$, gives every pair of lower and upper limits, and $\alpha=\beta=\delta$ gives every limit.
\Cref{thm:program,thm:pairs,thm:limitsets} are therefore not new as existence statements. Their content
is the rate $O(N^{-1/3})$, uniform in $\delta$; the denominator bound $q_k\le pk$ with two explicit
targets in \Cref{thm:program,thm:pairs}; the bound $k$ on the period of the state in
\Cref{thm:limitsets}; and the explicit arithmetic schedule of targets.
\end{remark}

\section{Quadratic images}\label{sec:quadratic}

Throughout this section $p\nmid b$ is an odd prime, $t=\ord_p(b)$, $v_0=v_p(b^t-1)\ge1$ and $d\ge3$ is an integer.
For an integer $a\ge1$ put $w_a(i)=p^{-av_p(i)}$ and
\[
X_a=\sum_{i\ge1}w_a(i)\,b^{-ti^d},
\]
the number $X_a$ of \Cref{sec:switch} with $n_k=tk^d$. All the $X_a$ share the schedule $n_i=ti^d$. We
use the lifting of exponents in the form $v_p(b^{tm}-1)=v_0+v_p(m)$ for $m\ge1$ (\Cref{sec:prelim}).

\begin{theorem}[Quadratic images]\label{thm:quadratic}\evidence{Paper proof}
Let $S$ be a finite set of positive integers and let $P\in\Q[z_a:a\in S]$ be nonconstant of total
degree at most two. Then $P((X_a)_{a\in S})$ is normal to base $b$. Consequently:
\begin{enumerate}[label=\textup{(\roman*)},nosep]
\item every rational polynomial of degree one or two in $X_s$, $s\ge1$, is normal to base $b$;
\item the distinct nonconstant monomials of degree at most two in $(X_a)_{a\in S}$ are jointly normal:
the points $(\fr{b^nm_1},\ldots,\fr{b^nm_r})$, $n\ge0$, are uniformly distributed in $\T^r$ for the
values $m_1,\ldots,m_r$ of these monomials;
\item an algebraic $X_s$ would have degree at least three.
\end{enumerate}
\end{theorem}

\begin{corollary}\label{cor:flagship}\evidence{Paper proof}
The number $\xi^2$, with $\xi=\sum_{k\ge1}3^{-v_3(k)}2^{-2k^3}$, is normal to base $2$, and the binary
shifts of $(\xi,\xi^2)$ are uniformly distributed in $\T^2$.
\end{corollary}

\begin{proof}
Here $b=2$, $p=3$, $t=2$, $d=3$ and $\xi=X_1$; apply \Cref{thm:quadratic}(i) and (ii).
\end{proof}

\paragraph{Comparison with exponential schedules.} Mira Cult's note \cite[Theorem~1.1]{MiraPoly}
(a conventional proof, not refereed or formally verified) proves that every nonconstant rational
polynomial in $x=\sum_{k\ge1}q^{-k}b^{-c^k}$ is normal to base $b$ when $q>2$ and some prime divides $c$
and $q$ but not $b$; this includes the Stoneham numbers $\alpha_{b,c}$ with $c\ge3$, $\gcd(b,c)=1$.
There the schedule $c^k$ is exponential and the weight exponents increase with $k$, and the leading
$p$-adic part of a truncated power of $x$ is controlled by a nonintegral algebraic tail, which does not
vanish. On the polynomial schedule $ti^d$ with valuation weights the analogous leading residue of
$R_n(\xi^2)$ vanishes on a set of shifts of positive lower density (\Cref{sec:intro}), so no single
surviving pole is available; the proof below replaces it by the measure-theoretic statement that the
limiting mass has null level sets (\Cref{prop:nulllevel}). Neither result contains the other.

\begin{remark}[An earlier bound]\label{rem:history}
Before the proof of \Cref{thm:quadratic} was found, an exact enumeration of the cells of $u$ on which
the normalized prefix of $\xi^2$ (see \Cref{lem:freeze} below) vanishes modulo $3^{21}$ gave the bound
$\limsup_ND_N(\xi^2;2)<2.94\cdot10^{-10}$. \Cref{cor:flagship} supersedes that
bound. The theorem gives no rate.
\end{remark}

The proof occupies the rest of this section. Products of two series are expanded into ordered pairs of
indices, and we write $R_n(Z)$ for the shifted prefix of a series $Z=\sum_ma_mb^{-m}$, as in
\Cref{prop:criterion}. For a product the coefficient at $b^{-m}$ is the sum of the products of
coefficients over the ordered pairs $(i,j)$ with $n_i+n_j=m$.

\subsection{Normalization and the limiting mass}

For $n\ge t$ let
\[
K(n)=\lfloor(n/t)^{1/d}\rfloor,\qquad J=J(n)=\lfloor\log_pK(n)\rfloor,\qquad u=u(n)=\frac{n}{tp^{dJ}} .
\]
Then $p^J\le K(n)<p^{J+1}$, so $u\in[1,p^d)$. Let $\Gamma$ be the set of positive rationals $a/p^h$ with
$h\ge0$ and $p\nmid a$, and call $\lambda(a/p^h)=h$ the \emph{loss}. For integers $a,c,D\ge1$ define, for
$u\in[1,p^d)$,
\[
M_{a,c}(u)=\sum_{\substack{(z,y)\in\Gamma^2\\ z^d+y^d\le u}}p^{a\lambda(z)+c\lambda(y)},\qquad
M^{\mathrm{lin}}_D(u)=\sum_{\substack{z\in\Gamma\\ z^d\le u}}p^{D\lambda(z)} .
\]
The conditions force $z,y<p$, so for each $L$ only finitely many terms are nonzero modulo $p^L$. The
sums therefore converge in $\Z_p$, and each reduction modulo $p^L$ is a right-continuous step function
of $u$ with finitely many jumps.

\begin{lemma}[Freezing]\label{lem:freeze}\evidence{Paper proof}
Let $a,c,D,L\ge1$ and let $n$ satisfy $J(n)\ge L+L/d$. Then
\[
p^{(a+c)J}b^{-n}R_n(X_aX_c)\equiv M_{a,c}(u)\pmod{p^L},\qquad
p^{DJ}b^{-n}R_n(X_D)\equiv M^{\mathrm{lin}}_D(u)\pmod{p^L}.
\]
Moreover $R_n(X_a)$ and $R_n(X_aX_c)$ are integral at every prime other than $p$, and
$v_p(R_n(X_a))\ge-aJ$, $v_p(R_n(X_aX_c))\ge-(a+c)J$.
\end{lemma}

\begin{proof}
By definition $b^{-n}R_n(X_aX_c)=\sum w_a(i)w_c(j)b^{-ti^d-tj^d}$ over the pairs with $t(i^d+j^d)\le n$.
Such indices satisfy $i,j\le K(n)<p^{J+1}$, so $\lambda_i=J-v_p(i)\ge0$ and
$p^{(a+c)J}w_a(i)w_c(j)=p^{a\lambda_i+c\lambda_j}$. This proves the integrality and the valuation bounds,
since $b^{n-n_i-n_j}$ is an integer. Write $z=i/p^J$, $y=j/p^J$; then $z,y\in\Gamma$ with losses
$\lambda_i,\lambda_j$, and $t(i^d+j^d)\le n$ if and only if $z^d+y^d\le u$. A term is nonzero modulo $p^L$
only if $a\lambda_i+c\lambda_j<L$, so $\lambda_i,\lambda_j<L\le J$; conversely every pair in $\Gamma^2$ with
losses less than $L$ arises from integers $i=p^Jz$, $j=p^Jy$. For such terms
$v_p(b^{ti^d}-1)=v_0+d(J-\lambda_i)\ge dL/d=L$, so $b^{-ti^d}\equiv1\pmod{p^L}$, and likewise for $j$.
The linear case is the same with one index.
\end{proof}

\subsection{Isolated jumps}

Fix an integer $D\ge2$ and integers $A_a$, $1\le a\le D/2$, not all zero, and $A_{\mathrm{lin}}$. Let
\[
M(u)=\sum_{1\le a\le D/2}A_aM_{a,D-a}(u)+A_{\mathrm{lin}}M^{\mathrm{lin}}_D(u).
\]
Let $a_*$ be the least $a$ with $A_a\ne0$ and $g=v_p(A_{a_*})$. It is convenient to use the variable
$x=(u-1)^{1/d}\in(0,X_{\max})$, $X_{\max}=(p^d-1)^{1/d}<p$; we write $M(x)$ for $M(1+x^d)$. A pair
$(z,y)\in\Gamma^2$ with $z^d+y^d>1$ contributes to $M(x)$ exactly for $x\ge x_{z,y}=(z^d+y^d-1)^{1/d}$, the
\emph{endpoint} of the pair; pairs with $z^d+y^d\le1$ contribute for all $x$. Similarly a linear term
$z>1$ contributes for $x\ge(z^d-1)^{1/d}$. A pair with $z=1$ has endpoint $y$.

For $H\ge1$ put $P_H=a_*H+g+1$. A \emph{primitive centre of level $H$} is a point $x_0=c/p^H\in(0,X_{\max})$
with $p\nmid c$.

\begin{lemma}[Primitive jumps]\label{lem:primjump}\evidence{Paper proof}
Let $H>g$. Modulo $p^{P_H}$, every jump of $M$ comes from one of the following: \textup{($\alpha$)} a pair
with a coordinate equal to $1$, whose other coordinate has loss at most $H$; \textup{($\beta$)} a pair
$(z,y)$ with $z\ne1\ne y$ and $\lambda(z)+\lambda(y)\le H$; \textup{($\gamma$)} a linear term with
$\lambda(z)\le H/2+g/D$. At a primitive centre $x_0$ of level $H$ the jumps of type $(\alpha)$ add up to
$A_{a_*}(p^{a_*H}+p^{(D-a_*)H})$, which is nonzero modulo $p^{P_H}$, and every other jump of type
$(\alpha)$ lies at distance at least $p^{-H}$ from $x_0$.
\end{lemma}

\begin{proof}
For the monomial of first order $a$, a pair of losses $(h,k)$ contributes $A_ap^{ah+(D-a)k}$. If $a=a_*$
this is nonzero modulo $p^{P_H}$ only if $a_*h+(D-a_*)k\le a_*H$, which forces $h+k\le H$ since
$D-a_*\ge a_*$. If $a>a_*$ and $A_a\ne0$, it is nonzero only if $ah+(D-a)k\le a_*H+g-v_p(A_a)$, which
forces $h+k\le(a_*H+g)/a<H$ for $H>g$. A linear term contributes $A_{\mathrm{lin}}p^{D\lambda(z)}$, nonzero
only if $D\lambda(z)\le a_*H+g$, and $a_*\le D/2$. Endpoints of pairs with a coordinate $1$ are the other
coordinates $y$, which lie in $p^{-H}\Z$ when $\lambda(y)\le H$; every point of $p^{-H}\Z$ other than
$x_0$ is at distance at least $p^{-H}$ from $x_0$, and the pair $(1,1)$ has endpoint $1\ne x_0$. At $x_0$
itself the pairs $(1,x_0)$ and $(x_0,1)$ are distinct; for the monomial $a_*$ they contribute
$A_{a_*}(p^{(D-a_*)H}+p^{a_*H})$, of valuation exactly $g+a_*H$ (if $a_*=D-a_*$ this is
$2A_{a_*}p^{a_*H}$, and $p$ is odd), and for $a>a_*$ their contributions have valuation at least
$v_p(A_a)+aH>a_*H+g$.
\end{proof}

Fix $\epsilon\in(0,1/4)$ and put $\delta_H=\epsilon/(Hp^H)$. A primitive centre $x_0$ of level $H$ is
\emph{isolated} if no endpoint of type $(\beta)$ or $(\gamma)$ lies within distance $\delta_H$ of $x_0$.
By \Cref{lem:primjump}, at an isolated centre $M\bmod p^{P_H}$ is constant on $(x_0-\delta_H,x_0)$ and on
$[x_0,x_0+\delta_H)$, and the two constants differ.

The equidistribution estimate behind the counting uses the third-derivative test.

\begin{lemma}[Third derivatives]\label{lem:third}\evidence{Paper proof}
There is an absolute constant $C_1$ with the following property. Let $\mathcal I$ be a set of at most $M$
consecutive integers, $M\ge1$, and let $g$ be real with three continuous derivatives on its convex hull,
$\lambda\le|g'''|\le2\lambda$ there for some $\lambda>0$. Then
$|\sum_{m\in\mathcal I}e(g(m))|\le C_1(M\lambda^{1/6}+M^{1/2}\lambda^{-1/6})$.
\end{lemma}

\begin{proof}
If $\lambda\ge1$ or $M\le\lambda^{-1/3}$ the bound is trivial, since then $M\le M\lambda^{1/6}$ or
$M\le M^{1/2}\lambda^{-1/6}$. Otherwise put $Q=\lfloor\lambda^{-1/3}\rfloor$, so $1\le Q\le M$. The
Weyl--van der Corput inequality \cite[Lemma~2.5 and (2.3.2), p.~10]{GK91} bounds $|\sum e(g(m))|^2$ by
$\frac{2M}{Q}\sum_{|q|<Q}|\sum_me(g(m+q)-g(m))|$, the inner sum over $m,m+q\in\mathcal I$. For $q\ne0$ the
function $g(m+q)-g(m)$ has second derivative $\int_m^{m+q}g'''$ of constant sign and modulus between
$|q|\lambda$ and $2|q|\lambda$, so the second-derivative test \cite[Theorem~2.2, p.~8]{GK91} bounds the inner sum
by $C(M(|q|\lambda)^{1/2}+(|q|\lambda)^{-1/2})$. Summing over $q$ gives
$|\sum e(g(m))|^2\ll M^2/Q+M^2(Q\lambda)^{1/2}+M(Q\lambda)^{-1/2}\ll M^2\lambda^{1/3}+M\lambda^{-1/3}$.
\end{proof}

\begin{lemma}[Isolated centres]\label{lem:isolated}\evidence{Paper proof}
Let $\mathcal K\subset(0,X_{\max})$ be a compact interval. There are $\epsilon\in(0,1/4)$ and $c_0>0$,
depending on $p,d,\mathcal K$, such that for every interval $I\subseteq\mathcal K$ of positive length and
every $H\ge H_0(I,M)$, at least $c_0|I|p^H$ isolated primitive centres of level $H$ lie in $I$.
\end{lemma}

\begin{proof}
There are at least $(1-1/p)|I|p^H-2$ primitive centres of level $H$ in $I$. Since $\delta_H<p^{-H}/2$,
each endpoint is within $\delta_H$ of at most one point of $p^{-H}\Z$, so it spoils at most one centre.
Endpoints of type $(\gamma)$ number $O(p^{H/2+g/D+1})=o(p^H)$. Let $I^+$ be the
$\delta_H$-neighbourhood of $I$, and fix a compact interval $\mathcal K'\subset(0,X_{\max})$ whose
interior contains $\mathcal K$, so that $I^+\subseteq\mathcal K'$ for large $H$. Let $B_H(I)$ be the number of pairs $(z,y)$ of type $(\beta)$
whose endpoint lies in $I^+$ and satisfies $\dist{p^Hx_{z,y}}\le\epsilon/H$, that is, lies within
$\delta_H$ of $p^{-H}\Z$. We show
\begin{equation}\label{eq:collision}
B_H(I)\le C_2\,\epsilon|I|p^H+o(p^H),
\end{equation}
with $C_2$ depending on $p,d,\mathcal K$ and the $o(p^H)$ uniform in $I\subseteq\mathcal K$. Choosing
$\epsilon=\min(1/8,(1-1/p)/(2C_2))$ then leaves at least $\frac12(1-1/p)|I|p^H-o(p^H)$ isolated centres,
which proves the lemma with $c_0=(1-1/p)/4$.

Write each pair as $(z,y)$ with $v=\lambda(z)\ge k=\lambda(y)$, counting both orders, so that $y$ is the
coarser coordinate; we let $z$ range over all of $p^{-v}\Z\cap(0,p)$, which only enlarges the count. For
fixed $y$ the endpoint is $f_y(z)=(z^d+B)^{1/d}$ with $B=y^d-1$; we call the set of $z$ for fixed $y$ a
\emph{row}. The rows of loss $k$ number at most $p^{k+1}$, and a row of loss $v$ has at most $p^{v+1}$
points. We use
\[
\sum_{\substack{k\le v\\v+k\le H}}p^{v+k}\ll_pHp^H,\qquad
\sum_{\substack{k\le v\\v+k\le H}}p^{v}\ll_pp^H,\qquad
\sum_{\substack{k\le v\\v+k\le H}}p^{k}\ll_pH^2p^{H/2}.
\]

\emph{Removals.} Charge every removed point as a collision. The rows with $v<2H/5$ contain
$O(H^2p^{4H/5})$ points in all. Put $\eta=H^{-6}$. Since $|y^d-1|\ge|y-1|$ for $y>0$, rows with
$|B|\le\eta$ have $0<|a-p^k|\le\eta p^k$ for $y=a/p^k$, hence there are at most $2\eta p^k$ of them for
each $k$, and they contain $O(\eta Hp^H)$ points. In each remaining row remove the points with
$z\le\eta$ or $|(d-2)B-(d+1)z^d|\le\eta^d$. Since $w^d-z^d\ge(w-z)^d$ for $w\ge z\ge0$, these conditions
define at most two intervals of total length $O_d(\eta)$, which contain $O(\eta p^v+1)$ grid points. The
total removal is $O(H^2p^{4H/5}+\eta Hp^H+H^2p^{H/2})=o(p^H)$.

\emph{Derivatives.} For $z>0$ with $f_y(z)>0$,
\[
f_y'=z^{d-1}f_y^{1-d},\qquad f_y''=(d-1)Bz^{d-2}f_y^{1-2d},\qquad
f_y'''=(d-1)Bz^{d-3}f_y^{1-3d}\bigl((d-2)B-(d+1)z^d\bigr).
\]
On the retained points with $f_y(z)\in\mathcal K'$ we have $z<p$, $|B|\le p^d$ and $f_y$ bounded above
and below, so $c_{\mathcal K}\eta^{2d-2}\le|f_y'''|\le C_{\mathcal K}$, using $d\ge3$. Moreover
$f_y''''=(d-1)Bz^{d-4}f_y^{1-4d}\,\Pi(z)$ with $\Pi$ a polynomial of degree $2d$ whose leading coefficient
is $(d+1)(d+2)\ne0$. Since $f_y$ is increasing, $\{z:f_y(z)\in I^+\}$ is an interval, and the retained
part of it splits into $O_d(1)$ intervals on which $f_y'''$ has constant sign and is monotone, hence into
$O(\log H)$ \emph{pieces} on which $\Lambda\le|f_y'''|\le2\Lambda$ for some $\Lambda\ge c_{\mathcal K}\eta^{2d-2}$.

\emph{Equidistribution in a piece.} Put $T=p^v$ and, for an integer $1\le l\le H^4$,
$g_l(m)=lp^Hf_y(m/T)$ on the integers $m$ of a piece, at most $pT$ of them. Then $\lambda_l\le|g_l'''|\le
2\lambda_l$ with $\lambda_l=lp^H\Lambda/T^3$. By \Cref{lem:third},
\begin{gather*}
\Bigl|\sum_me(g_l(m))\Bigr|\ll_pT\bigl(\lambda_l^{1/6}+T^{-1/2}\lambda_l^{-1/6}\bigr),\\
\lambda_l^{1/6}\ll H^{O(1)}p^{H/6-v/2}\le H^{O(1)}p^{-H/30},\qquad
T^{-1/2}\lambda_l^{-1/6}=(lp^H\Lambda)^{-1/6}\le H^{O(1)}p^{-H/6},
\end{gather*}
using $v\ge2H/5$, $l\le H^4$ and $\Lambda\ge H^{-O(1)}$. The Erdős--Turán inequality
\cite[Chapter~2, Theorem~2.5, p.~112]{KN} with $H^4$ frequencies bounds the number of $m$ in the piece
with $\dist{p^Hf_y(m/T)}\le\epsilon/H$ by
\[
\frac{2\epsilon}H\#\{m\}+O\Bigl(\frac{T}{H^4}+\sum_{l\le H^4}\frac1l\Bigl|\sum_me(g_l(m))\Bigr|\Bigr)
=\frac{2\epsilon}H\#\{m\}+O(TH^{-4}),
\]
since the arc $\{\theta\in\T:\dist\theta\le\epsilon/H\}$ has length $2\epsilon/H$. Summed over the
$O(\log H)$ pieces of a row, the error is $O(TH^{-3})$, and summed over the rows it is
$O(H^{-3}\sum p^{k+v+1})=O(p^H/H^2)$.

\emph{Main term.} For fixed losses $(v,k)$ the number of all grid points
$(z,y)\in(p^{-v}\Z\times p^{-k}\Z)\cap(0,p)^2$ with endpoint in $I^+=[x_1,x_2]$ is the number of lattice
points with $1+x_1^d<z^d+y^d\le1+x_2^d$, up to the endpoint convention. The area of
$\{z,y>0:z^d+y^d\le r^d\}$ is a constant times $r^2$, so the region has area
$O((1+x_2^d)^{2/d}-(1+x_1^d)^{2/d})=O_{\mathcal K}(|I^+|)$. Each sublevel set lies under a decreasing
graph, so counting by columns and comparing monotone sums with integrals gives
$p^{v+k}\cdot\mathrm{area}+O(p^{v+1}+p^{k+1})$ for its lattice count. Hence the main terms sum to at most
\[
\frac{4\epsilon}H\sum_{\substack{k\le v\\v+k\le H}}\Bigl(C|I^+|p^{v+k}+O(p^v+p^k)\Bigr)
\le C_2\epsilon|I|p^H+O(\epsilon p^H/H)+O(\delta_Hp^H).
\]
Together with the removals and the error terms this proves \eqref{eq:collision}.
\end{proof}

\subsection{Null level sets}

\begin{proposition}[Null level sets]\label{prop:nulllevel}\evidence{Paper proof}
For every $w\in\Z_p$ the set $\{u\in[1,p^d):M(u)=w\}$ has Lebesgue measure zero.
\end{proposition}

\begin{proof}
Fix a compact interval $\mathcal K\subset(0,X_{\max})$, let $\epsilon$ and $c_0$ be as in
\Cref{lem:isolated}, and let $G_H$ be the union of the intervals $(x_0-\delta_H,x_0+\delta_H)$ over the
isolated primitive centres $x_0\in\mathcal K$ of level $H$.

\emph{Almost every point of $\mathcal K$ lies in infinitely many $G_H$.} Let $I\subseteq\mathcal K$ be an
interval. The intervals about isolated centres in the middle half of $I$ are disjoint and contained in
$I$ for large $H$, so by \Cref{lem:isolated}
\[
|I\cap G_H|\ge\tfrac12c_0|I|p^H\cdot2\delta_H=c_0\epsilon|I|/H\qquad(H\ge H_1(I)).
\]
Let $N_H$ be the $\delta_H$-neighbourhood of all of $p^{-H}\Z$. For $n<m$ with $|I|p^n\ge1$, at most
$|I|p^n+2\le3|I|p^n$ intervals of $N_n$ meet $I$, and each meets at most $2\delta_np^m+2$ intervals of
$N_m$, of length $2\delta_m$; hence
\[
|I\cap G_n\cap G_m|\le|I\cap N_n\cap N_m|\le3|I|p^n(2\delta_np^m+2)2\delta_m
\le12|I|\Bigl(\frac{\epsilon^2}{nm}+\frac{\epsilon p^{n-m}}m\Bigr),
\]
and $|I\cap G_n|\le6\epsilon|I|/n$. For $L\le H\le U$ with $L\ge H_1(I)$ put $S=\sum_{H=L}^U\mathbf1_{G_H}$
and $A=\sum_{H=L}^U1/H$. Then $\int_IS\ge c_0\epsilon|I|A$ and, using $\sum_{n<m}p^{n-m}/m\le A/(p-1)$,
$\int_IS^2\le|I|(12\epsilon^2A^2+C\epsilon A)$. By the Cauchy--Schwarz inequality
$|I\cap\{S>0\}|\ge(\int_IS)^2/\int_IS^2$, and letting $U\to\infty$, so that $A\to\infty$,
\[
\Bigl|I\cap\bigcup_{H\ge L}G_H\Bigr|\ge\frac{c_0^2}{12}|I| .
\]
The left side decreases in $L$, so this holds for every $L$, with a constant independent of $I$ and $L$.
Hence the set $\mathcal K\setminus\bigcup_{H\ge L}G_H$ has density at most $1-c_0^2/12$ in every interval
of $\mathcal K$, so it has no density point in the interior of $\mathcal K$, and by the Lebesgue density
theorem it is null. Taking the union over $L$ proves the claim.

\emph{Level sets.} The set $\Lambda_w=\{x\in\mathcal K:M(x)=w\}$ is the intersection over $L$ of the
finite unions of intervals $\{M\equiv w\bmod p^L\}$, hence Borel. Let $x\in\Lambda_w$ lie in $G_H$, within
$\delta_H$ of an isolated centre $x_0$ with $H>g$. The interval of radius $2\delta_H$ about $x$ contains
$(x_0-\delta_H,x_0+\delta_H)$. On the two halves of this interval $M\bmod p^{P_H}$ takes two different
constant values, so one whole half, of length $\delta_H$, misses $\Lambda_w$, and $\Lambda_w$ has relative
measure at most $3/4$ in the interval of radius $2\delta_H$ about $x$. If $x$ lies in infinitely many
$G_H$, it is therefore not a density point of $\Lambda_w$. Since almost every point of $\Lambda_w$ is a
density point and almost every point of $\mathcal K$ lies in infinitely many $G_H$, $\Lambda_w$ is null.
Exhausting $(0,X_{\max})$ by countably many compact intervals and returning to $u=1+x^d$, a $C^1$
diffeomorphism onto $(1,p^d)$, proves the proposition.
\end{proof}

\subsection{Proof of the theorem}

\begin{proof}[Proof of \Cref{thm:quadratic}]
\emph{Reduction.} Let $P$ have rational coefficients with common denominator $Q_0$, and write
$Q_0=Q_bQ_1$ with $\gcd(Q_1,b)=1$ and $Q_b\mid b^e$ for some $e$. Then $b^eP=Y/Q_1$ with $Y$ an integral
polynomial in the $X_a$, and $b^eP$ is normal if and only if $P$ is, since multiplying by $b^e$ shifts the
orbit by $e$ steps. If $r=\alpha/q$ in lowest terms, then $r/Q_1$ has reduced denominator between $q$ and
$qQ_1$, dividing $qQ_1$. Hence it suffices to verify the hypotheses of \Cref{prop:criterion} for the
series of $Y$, with $c=p$, and then for $Y/Q_1$ with $c=pQ_1$. We therefore assume that $P$ has integer
coefficients and verify (i)--(iii) for $Y=P((X_a))$.

\emph{Hypotheses (i) and (ii).} The positions of nonzero coefficients of $Y$ through $N$ lie in
$\{0\}\cup\{ti^d\le N\}\cup\{t(i^d+j^d)\le N\}$, which has $O(N^{2/d})$ elements; since $d\ge3$, (i) holds
with $\eta=1-2/d$. A coefficient at $b^{-m}$ is bounded by the sum of $|P|$'s coefficients times
$1+\#\{(i,j):t(i^d+j^d)=m\}\le1+m$. By \Cref{lem:freeze}, $q_n$ divides $p^{DJ}$, where $D$ is the largest
weight defined below, and $p^{DJ}\le K(n)^D\le(n+1)^D$.

\emph{Weights.} Give $X_a$ weight $a$ and $X_aX_c$ weight $a+c$, and let $D$ be the largest weight of a
monomial with nonzero coefficient in $P$. By \Cref{lem:freeze} every monomial of weight at most $D-1$
contributes to $p^{DJ}b^{-n}R_n(Y)$ a term of valuation at least $J$.

\emph{Case 1: the only monomial of weight $D$ is linear,} $A_{\mathrm{lin}}X_D$. By \Cref{thm:den}, $R_n(X_D)$
has $p$-adic valuation exactly $-DJ$, so $v_p(R_n(Y))=v_p(A_{\mathrm{lin}})-DJ$ for all large $n$. Hence
$q_n\to\infty$ and (iii) holds.

\emph{Case 2: some monomial $X_aX_{D-a}$ has a nonzero coefficient $A_a$} ($1\le a\le D/2$). Let $M$ be
the mass of the preceding subsection built from these coefficients and the coefficient
$A_{\mathrm{lin}}$ of $X_D$ (zero if absent). Fix $M_1\ge1$ and $L\ge1$. For $n$ with $J\ge L(1+1/d)$,
\Cref{lem:freeze} and the preceding remark give
\[
p^{DJ}b^{-n}R_n(Y)\equiv M(u(n))\pmod{p^{L}} .
\]
If $M(u(n))\not\equiv0\pmod{p^{L}}$, the left side has valuation at most $L-1$, so
$q_n\ge p^{DJ-L+1}$, which exceeds $M_1$ once $J$ is large. Therefore, for all large $n$, $q_n<M_1$ implies
$u(n)\in E_L=\{u:M(u)\equiv0\pmod{p^L}\}$. The set $E_L$ is a finite union of, say, $\kappa_L$ intervals.
The integers $n$ with $J(n)=j$ form the interval $[tp^{dj},tp^{d(j+1)})$, and those among them with
$u(n)\in E_L$ number at most $tp^{dj}|E_L|+\kappa_L$. If $tp^{dJ_0}\le N<tp^{d(J_0+1)}$, summing over
$j\le J_0$ gives
\[
\#\{n<N:q_n<M_1\}\le O_{L,M_1}(1)+\frac{tp^{d(J_0+1)}}{p^d-1}|E_L|+\kappa_L(J_0+1)
\le\Bigl(\frac{p^d}{p^d-1}|E_L|+o(1)\Bigr)N .
\]
The sets $E_L$ decrease to $\{M=0\}$, which is null by \Cref{prop:nulllevel}, so $|E_L|\to0$, and (iii)
follows. \Cref{prop:criterion} proves normality.

\emph{Consequences.} For (ii), every nonzero integer combination of distinct nonconstant monomials of
degree at most two is a nonconstant polynomial of total degree at most two, hence normal; uniform
distribution of the vector of shifts follows from Weyl's criterion in several dimensions \cite[Chapter~1,
Theorem~6.2, p.~49]{KN}. For (iii), if $X_s$ were algebraic of degree at most two, a rational polynomial
of degree at most two would vanish at $X_s$, and $0$ is not normal.
\end{proof}

\begin{remark}[Scope]\label{rem:quadscope}
The argument gives no rate and no information on the zero set of $M$ beyond its measure. For an
$r$-fold product the number of grid tuples of total loss at most $H$ is of order $H^{r-1}p^H$, and
isolating a single primitive grid at radius $1/(H^{r-1}p^H)$ gives a convergent sum of neighbourhood
measures when $r\ge3$, so the second-moment step fails; when $p\mid r$ the basic jump $rp^{sH}$ also loses
a factor $p$. Normality of $\xi^3$ is open in the sources we examined.
\end{remark}

\section{Infinite sums}\label{sec:sums}

Let $s_3=2$ and $s_\ell=1$ for primes $\ell\ge5$, and for every odd prime $\ell$ put
\[
Y_\ell=\sum_{j\ge1}\frac{1}{\ell^{s_\ell j}\,2^{16^{\ell^j}}},\qquad
Y=\sum_{\ell\ \text{odd prime}}Y_\ell .
\]
The insertions are indexed by the odd prime powers $m=\ell^j$ (the \emph{labels}); the label $m$ sits
at position $16^m$, the positions are distinct, and every coefficient is at most $1/5$. For $n\ge1$ let
$L=L(n)=\lfloor\log_{16}n\rfloor$, so that the terms inserted by shift $n$ are those with label at most
$L$, and for $\ell\le L$ let $J_\ell=\lfloor\log_\ell L\rfloor$.

\begin{lemma}\label{lem:Yden}\evidence{Paper proof}
Let $r_\ell$ $(\ell$ odd prime$)$ be finitely supported rationals with odd denominators, and let $R_n$ be
$2^n$ times the sum of the terms of $\sum_\ell r_\ell Y_\ell$ with label at most $L(n)$. For every $\ell$
with $r_\ell\ne0$, $v_\ell(R_n)=v_\ell(r_\ell)-s_\ell J_\ell$ for all sufficiently large $n$, and $R_n$ is
integral at $2$ and at every prime outside the support of $(r_\ell)$ and the denominators of the
$r_\ell$. With $r_\ell=1$ for all $\ell$, the reduced denominator of $R_n$ is exactly
$q(n)=3^{2J_3}\prod_{5\le\ell\le L}\ell^{J_\ell}$ for every $n$ with $L(n)\ge3$.
\end{lemma}

\begin{proof}
The terms $r_\ell\ell^{-s_\ell j}2^{n-16^{\ell^j}}$ with $\ell^j\le L$ are $2$-integral. At a prime $\ell$,
the terms from other axes have valuation at least $\min_{\ell'}v_\ell(r_{\ell'})$, a constant; on the axis
$\ell$ the term with $j=J_\ell$ has valuation $v_\ell(r_\ell)-s_\ell J_\ell$ and the others have strictly
larger valuation. Since $J_\ell\to\infty$, the top term has the unique least valuation for large $n$.
When all $r_\ell=1$ the other axes are $\ell$-integral and the conclusion holds as soon as $\ell\le L$.
\end{proof}

We use the elementary bound $\prod_{\ell\le x}\ell<4^x$ over primes: the primes in $(m+1,2m+1]$ divide
$\binom{2m+1}{m}\le4^m$, and induction on $x$ completes the proof. It gives
$q(n)\le3^{J_3}\prod_{\ell\le L}\ell^{J_\ell}\le L^{\sqrt L+1}4^L$, since only primes $\ell\le\sqrt L$
have $J_\ell\ge2$.

\begin{theorem}[All odd primes]\label{thm:Y}\evidence{Lean-supported paper proof}
The number $Y$ is normal to base $2$ and $D_N(Y;2)\ll1/\log(2N)$ for all $N\ge1$. Moreover $Y$ is a
Liouville number.
\end{theorem}

\begin{proof}
Let $L(n)\ge3$, write $q=q(n)=3^eQ$ with $e=2J_3$ and $Q$ odd and prime to $3$, and let
$\tau=\lcm(2,\ord_Q(2))$, with the conventions $\ord_1(2)=1$ and that an empty maximum is zero (both
apply when $L\le4$, where $Q=1$). Then $2^\tau\equiv1\pmod{3Q}$ and, lifting the exponent from $v_3(4-1)=1$,
\[
c=v_3(2^\tau-1)=1+v_3(\ord_Q(2))\le1+\max_{5\le\ell\le L}v_3(\ell-1)\le1+J_3,
\]
because the order of $2$ modulo $\ell^f$ divides $\ell^{f-1}(\ell-1)$ and $3^{v_3(\ell-1)}\le\ell-1<L$.
Hence $e-c\ge J_3-1\ge0$, and \Cref{lem:grid} shows that any $L'$ consecutive points $\fr{a2^j/q}$
with $3\nmid a$ (in particular with $a$ prime to $q$, as for the reduced prefixes) have count error at most $L'3^{1-J_3}+\tau3^{e-c}\le9L'/L+q$ on every half-open arc, using
$3^{J_3+1}>L$ and $\tau\le2Q$. The doubled order at $3$ is what makes the grid grow: the loss $c$ can be
as large as $J_3$, but $e=2J_3$. The Lean-supported steps are the order and count statements cited in
\Cref{lem:grid}.

Let $N$ be large, $N_0=\lceil\sqrt N\rceil$, and apply \Cref{lem:blocks}: leave $[0,N_0)$ uncovered and
split $[N_0,N)$ at the insertion positions. Fewer than $\log_{16}(2N)$ labels have position below $2N$,
so there are $O(\log N)$ blocks; on each, $R_{n+1}=2R_n$ and the models are consecutive points
$\fr{a2^j/q}$ with $q$ constant, and $L(n)\ge\frac12\log_{16}N-1$. With $w=\lceil\log_2N\rceil+1$, a
shift $n<N$ with no insertion in $(n,n+w]$ satisfies $0\le2^nY-R_n<1/N$, and at most $w\log_{16}(2N)$
shifts fail this. Since $L(N)\le\log_{16}N$, $q\le L^{\sqrt L+1}4^L\le N^{1/2+o(1)}$, so
\[
D_N(Y)\le\frac{N_0+w\log_{16}(2N)+O(N/\log N)+O(\log N)N^{1/2+o(1)}}{N}+\frac2N\ll\frac1{\log N}.
\]
Small $N$ are absorbed in the constant.

For the Liouville property, given $h\ge1$ choose distinct primes $r_1,\ldots,r_{2h}$ and, by the Chinese
remainder theorem, arbitrarily large $z$ with $r_{2i-1}r_{2i}\mid z+i$ for $1\le i\le h$. None of
$z+1,\ldots,z+h$ is a prime power, so there are consecutive labels $u<v$ with $v-u>h$ and $u$
arbitrarily large. The truncation $P$ through label $u$ has denominator at most $2^{16^u}q$ with
$\log_2q=O(u)$, while $0<Y-P\le2^{1-16^v}\le2^{1-16^{h}\cdot16^u}$. As $h$ is arbitrary and $Y$ is
irrational, $Y$ is a Liouville number.
\end{proof}

If only $3$ and the primes $\ell\equiv2\pmod3$ are used, $3\nmid\ord_Q(2)$, so $c=1$ and the weight
$3^{-j}$ already gives a growing grid at $3$; with positions $2^{\ell^{2j}}$ this yields the rate
$O((\log N)^{-1/2})$ by the same argument.

\begin{theorem}[Finite combinations]\label{thm:combos}\evidence{Lean-supported paper proof}
For every finite set $S$ of odd primes and rationals $r_\ell$ $(\ell\in S)$, not all zero, the number
$\sum_{\ell\in S}r_\ell Y_\ell$ is normal to base $2$. Consequently $(Y_\ell)_{\ell\in S}$ is jointly
normal: the points $(\fr{2^nY_\ell})_{\ell\in S}$ are uniformly distributed in $\T^S$.
\end{theorem}

\begin{proof}
Multiplying a number by $2^a$ changes $D_N$ by at most $a/N$, so we may assume that all $r_\ell$ have odd
denominators. Choose $\ell_0\in S$ with $r_{\ell_0}\ne0$, let $F$ be the finite set of odd primes dividing
$\prod_{\ell\in S}\ell$ or a denominator of some $r_\ell$, and let
$\beta=\max_{\ell\in F\setminus\{\ell_0\}}v_{\ell_0}(\ell-1)$ (zero if empty). By \Cref{lem:Yden}, for large
$n$ the reduced denominator of $R_n$ is $q=\ell_0^eQ$ with $e=s_{\ell_0}J_{\ell_0}-v_{\ell_0}(r_{\ell_0})\to\infty$
and $Q$ supported on $F\setminus\{\ell_0\}$. The order $\ord_Q(2)$ divides $\lcm_\ell\ell^{f_\ell-1}(\ell-1)$,
so its $\ell_0$-valuation is at most $\beta$. With $t_0=\ord_{\ell_0}(2)$ and $\tau=\lcm(t_0,\ord_Q(2))$,
lifting the exponent gives $c=v_{\ell_0}(2^\tau-1)=v_{\ell_0}(2^{t_0}-1)+v_{\ell_0}(\tau/t_0)\le
v_{\ell_0}(2^{t_0}-1)+\beta=:c^*$, independent of $n$. \Cref{lem:grid} gives count error at most
$L'\ell_0^{c^*-e}+\tau\ell_0^{e-c}\le L'\ell_0^{c^*-e}+q$, and $q\le(\log N)^{O(1)}$ for $n<N$. The block and tail
argument of \Cref{thm:Y}, with $O(\log N)$ insertions and the tail multiplied by $\sum|r_\ell|$, gives
$D_N\to0$. Joint normality follows by applying this to the integer combinations $\sum h_\ell Y_\ell$,
$h\ne0$, together with Weyl's criterion and its multidimensional form
\cite[Chapter~1, Theorems~2.1 and~6.2, pp.~7 and~49]{KN}.
\end{proof}

Only a fixed finite set of primes enters $Q$ here, so the loss $c$ is bounded. For the infinite sum
$Y$ the loss can grow, and the doubled order at $3$ compensates for it.

\begin{theorem}[Rational totals]\label{thm:rationaltotal}\evidence{Lean-supported paper proof}
For every real $R>0$ there are rationals $c_\ell>0$, indexed by the odd primes, with
$\sum_\ell c_\ell Y_\ell=R$. For rational $R$ the total is not normal, although every finite subfamily
of the positive terms $c_\ell Y_\ell$ is jointly normal.
\end{theorem}

\begin{proof}
For positive generators $g_i$ and $R>0$, choose inductively rationals with $e_{i-1}/2<c_ig_i<e_{i-1}$,
where $e_0=R$ and $e_i=e_{i-1}-c_ig_i$; then $0<e_i<R2^{-i}$. This is the kernel-verified
\lean{PositiveSummation}{15}{\ttt{hasSum\_positive\_rational\_multiples}}, applied to an enumeration of
the $Y_\ell$. Joint normality of finite subfamilies is \Cref{thm:combos}, since nonzero rational scaling
preserves its hypothesis. A rational number has an eventually periodic expansion and is not normal.
\end{proof}

The tails of a series of normal numbers cannot all be negligible. For real $z$ let
$M_N(z)=N^{-1}\sum_{n<N}\dist{b^nz}$ and $\overline M(z)=\limsup_NM_N(z)$.

\begin{proposition}[Consecutive tails]\label{prop:tails}\evidence{Lean-supported paper proof}
Let $\sum_iS_i$ be a convergent series of numbers normal to base $b$, with tails $E_k=\sum_{i>k}S_i$.
Then for every $k$,
\[
\overline M(E_k)+\overline M(E_{k+1})\ge\frac14\qquad\text{and}\qquad
\limsup_N\bigl(D_N(E_k)+D_N(E_{k+1})\bigr)\le1 .
\]
In particular at most one of any two consecutive tails $z$ satisfies $D_N(z)\to1$, that is, by
\Cref{thm:sharpcited}, $\sup_x|D_N(x+z)-D_N(x)|\to0$.
\end{proposition}

\begin{proof}
Since $S_{k+1}=E_k-E_{k+1}$, $\dist{b^nS_{k+1}}\le\dist{b^nE_k}+\dist{b^nE_{k+1}}$. By uniform distribution
\cite[Chapter~1, Corollary~1.1, p.~3]{KN} the averages of the left side tend to $\int_0^1\dist{y}\,dy=1/4$.
For the second inequality apply \Cref{thm:sharpcited} with $x=E_k$ and $z=-E_{k+1}$: since
$D_N(-E_{k+1})=D_N(E_{k+1})$, it gives $D_N(S_{k+1})\ge D_N(E_k)+D_N(E_{k+1})-1$, and $D_N(S_{k+1})\to0$.
The Lean-supported step is \Cref{thm:sharpcited}.
\end{proof}

In \Cref{thm:Y} normality of the total therefore comes from the arithmetic grid at $3$, which survives
every new prime, and not from smallness of the tails. The translations $z$ with $D_N(z)\to1$ are
classified in \cite[Theorem~3.3]{CompanionH}; in base $2$ they are exactly the $z$ with $M_N(z)\to0$.

\section{Questions}\label{sec:questions}

\begin{question}
Give an explicit rate for $D_N(\xi^2;2)$. Is $\xi^3$ normal to base $2$?
\end{question}

\begin{question}
Is $\xi$ transcendental? Is it a Liouville number? (Normality excludes only rationality, and
\Cref{thm:quadratic} excludes algebraic degree at most two.)
\end{question}

\begin{question}
For $s\ge1$ and $2\le d\le s+1$, what is the true order of $D_N(X_s)$ for $n_k=tk^d$? Is the exponent
$1/3$ in \Cref{thm:B,thm:C} optimal for these particular numbers?
\end{question}

\section{Formal verification and data availability}\label{sec:formal}

The Lean~4 \cite{Lean4} declarations cited above, built with Mathlib \cite{Mathlib}, are in the
repository \url{https://github.com/eldernos/MathematicsResearch}. Links to the modules
\ttt{RetirementDenominator}, \ttt{Phase18Synthetic}, \ttt{SharpTranslation}, \ttt{SmoothOrbit} and
\ttt{PositiveSummation} are pinned to the commit
\begin{center}\ttt{\leancommit},\end{center}
and links to \ttt{SurvivingGrid} and \ttt{PrimePowerGrid} to the later commit
\begin{center}\ttt{\phaseXIXcommit}.\end{center}
At both commits the full build completed and the axiom audit accepted only \ttt{propext},
\ttt{Classical.choice} and \ttt{Quot.sound}, and an independent kernel replay covered all modules. The
theorem verification map, the novelty ledger and the reproduction commands are in the files
\ttt{VERIFICATION.md} and \ttt{LEDGER.md} next to the source of this paper (directory
\ttt{papers/07-synthetic-expansions} of the repository).

The kernel-verified statements are
\Cref{prop:group,thm:den,lem:short,thm:switch,lem:program,thm:sharpcited}; the last is quoted from
\cite{CompanionH}. All rates (\Cref{thm:example,thm:rate,thm:program,thm:Y}), the limit-set theorems,
\Cref{prop:criterion} beyond the short-orbit estimate, the passage from repunits to reduced raw
coefficients after \Cref{prop:group}, \Cref{sec:quadratic} and \Cref{sec:sums} have written proofs. In particular \Cref{thm:quadratic} and \Cref{cor:flagship} are paper proofs and are not
formalized. Exact rational computations exercising these statements are
\repo{scripts/phase18_principal.py}{\ttt{scripts/phase18\_principal.py}} and
\repo{scripts/phase18_programming.py}{\ttt{scripts/phase18\_programming.py}} at the first commit, and
\href{https://github.com/eldernos/MathematicsResearch/tree/\phaseXIXcommit/scripts/phase19_grids.py}{\ttt{scripts/phase19\_grids.py}}
and
\href{https://github.com/eldernos/MathematicsResearch/tree/\phaseXIXcommit/scripts/phase19_nonlinear.py}{\ttt{scripts/phase19\_nonlinear.py}}
at the second; they are finite checks and illustrations, not proofs of the asymptotic statements. The
enclosures of $N^{1/3}D_N(\xi)$ quoted after \Cref{thm:example} are computed by
\ttt{checks/xi\_discrepancy.py} in the manuscript directory (standard library only).

\section{Use of AI tools}

This manuscript was produced in an AI-assisted research project. Large language models proposed
conjectures and proof strategies, wrote the proofs and the text, implemented the exact computations
and Lean~4 formalizations, searched the literature and reviewed drafts across several model
providers. The repository owner directed the project and is responsible for its release. The Lean
kernel checked the declarations listed in the verification map. The written proofs, literature
comparisons and priority statements have been checked only by AI models from different providers.
They have not been checked by an independent human expert. Cross-model review is not peer review.

\bibliographystyle{plain}
\bibliography{references}
\end{document}
